\section{Lines at \texorpdfstring{$60^\circ$ and $90^\circ$}{60 and 90 degrees}}\label{sec:2}
% Long structured steps should leave space at a page foot instead of stretching proof spacing.
\raggedbottom

We follow Section 2 of [CGSS], with the vector encoding of Section~1.
We use the graph names and Lean names of \path{cgss/notation.md}; in Lean,
\texttt{Hyp24} is in the namespace \texttt{LineSystem}. In what follows,
$V$ is a finite-dimensional real inner product space. We write $(\Zx,\Zy)$ for its inner
product, as the paper does, since angle brackets are taken by step numbers.
A \emph{line system} $S$ in $V$ is a finite set $\Phi\subseteq V$ such that
\[
 \Phi=-\Phi,\qquad (\Zr,\Zr)=2\quad(\Zr\in\Phi),\qquad
 (\Zr,\Zs)\in\{-2,-1,0,1,2\}\quad(\Zr,\Zs\in\Phi).
\]
\label{def:line-system}\lean{LineSystem}\uses{}
We count lines as antipodal pairs $\{\Zr,-\Zr\}$, so the number of lines is $|\Phi|/2$.
\lean{LineSystem.numLines}
We do not require $\Phi$ to span $V$ until Theorem 2.13.
We call $S$ \emph{indecomposable} if, for all $P,Q\subseteq V$ with $P\cup Q=\Phi$ and
$(\Zr,\Zs)=0$ for all $\Zr\in P$, $\Zs\in Q$, we have $P=\varnothing$ or $Q=\varnothing$.
\label{def:indecomposable}\lean{LineSystem.IsIndecomposable}\uses{def:line-system}
\new\ This intrinsic definition agrees with the paper's orthogonal-subspace definition
when $\Phi$ spans $V$; without spanning, the paper's literal wording allows the partition
into the whole system in its proper span and an empty part. We also require finiteness
explicitly: in finite dimension this loses no line systems, since their normalized signed
vectors are pairwise at distance at least 1 on the compact unit sphere.
An isomorphism of line systems is a bijection of their vector sets preserving inner products.
\label{def:line-iso}\lean{LineSystem.Iso}\uses{def:line-system}

\paragraph{Elementary identities.}\label{facts:inner}
We will use bilinearity, symmetry, nonnegativity of squared norm, and the implication
$(\Zx,\Zx)=0\Rightarrow\Zx=0$. These are, respectively, Mathlib's\par
\begingroup\raggedright
\texttt{inner\_add\_left}, \texttt{inner\_add\_right},
\texttt{inner\_sub\_left}, \texttt{inner\_sub\_right},
\texttt{inner\_neg\_left}, \texttt{inner\_neg\_right},
\texttt{real\_inner\_smul\_left}, \texttt{real\_inner\_smul\_right},
\texttt{real\_inner\_comm}, \texttt{real\_inner\_self\_nonneg}, and
\texttt{inner\_self\_eq\_zero}, all in
\path{Mathlib/Analysis/InnerProductSpace/Basic.lean}.\par\endgroup
Here \texttt{inner\_self\_eq\_zero} is the equivalence
$(\Zx,\Zx)=0\Longleftrightarrow\Zx=0$; its \texttt{NormedAddCommGroup}
hypothesis holds for our real inner product spaces.
In particular, for vectors of squared norm 2,
\begin{equation}\label{eq:endpoints}
 (\Zr,\Zs)=2\Longleftrightarrow\Zr=\Zs,
 \qquad (\Zr,\Zs)=-2\Longleftrightarrow\Zr=-\Zs.
\end{equation}
Indeed, $(\Zr-\Zs,\Zr-\Zs)=4-2(\Zr,\Zs)$ and
$(\Zr+\Zs,\Zr+\Zs)=4+2(\Zr,\Zs)$ by bilinearity; the forward implications follow
from zero squared norm, and the reverse implications follow by substitution.
Thus distinct antipodal pairs represent distinct lines with angles $60^\circ$ or $90^\circ$.
We will also use the \emph{finite-span lemma}: for a finite family $(\Zr_i)_{i\in I}$
in $V$, every member of its span is $\sum_{i\in I}\alpha(i)\Zr_i$ for some
$\alpha:I\to\re$. This is
\path{Submodule.mem_span_range_iff_exists_fun} in
\path{Mathlib/LinearAlgebra/Finsupp/LinearCombination.lean}.

\begin{cgss}{Definition}{2.1}\label{def:2.1}
 \lean{LineSystem.IsStarClosed}\uses{def:line-system}
 \assume{$S$ is a line system in a finite-dimensional real inner product space $V$.}
 We call $S$ \emph{star-closed} if
 \[
  \forall\Zr,\Zs\in\Phi,\quad (\Zr,\Zs)=-1\ \Longrightarrow\ \Zr+\Zs\in\Phi.
 \]
 The third line of the star on $\Zr$ and $\Zs$ is the line of $\Zr+\Zs$:
 its squared norm is $2+2-2=2$, and its inner product with each summand is 1.
 The three vectors lie in the plane spanned by $\Zr,\Zs$.
\end{cgss}

\begin{cgss}{Definition}{2.2}\label{def:2.2}
 \lean{LineSystem.IsMaximal}\uses{def:line-system}
 \assume{$S$ is a line system in a finite-dimensional real inner product space $V$.}
 We call $S$ \emph{maximal} if, for every $\Zv\in\operatorname{span}(\Phi)$,
 \[
 (\Zv,\Zv)=2\ \text{ and }\
 \bigl(\forall\Zr\in\Phi,\ (\Zv,\Zr)\in\{-2,-1,0,1,2\}\bigr)
 \quad\Longrightarrow\quad \Zv\in\Phi.
 \]
 We include $\pm2$ to allow a line already in $S$; by \eqref{eq:endpoints} these values
 occur only on that same line. \repair\ The paper takes the lines to span $\re^n$ (the first
 paragraph of Section 2) and asks for maximality in $\re^n$; since we do not require $\Phi$ to
 span $V$, we ask for maximality within $\operatorname{span}(\Phi)$. In the whole of $V$ the
 condition would be too strong: $\Ze_1+\Ze_2$ is at $60^\circ$ or $90^\circ$ to every line of
 $\typ{A}_\ell$ in $\re^{\ell+1}$, and Theorem 3.6 would fail (errata, Implicit hypothesis 7). Thus this is
 Definition 2.2 in the signed vector encoding.
\end{cgss}

\begin{cgss}{Lemma}{2.3}\label{lem:2.3}
 \lean{LineSystem.IsMaximal.isStarClosed}\uses{def:2.1,def:2.2,facts:inner}
 \assume{\begin{pfenum}
  \item $S$ is a line system in a finite-dimensional real inner product space $V$.
  \item $S$ is maximal.
 \end{pfenum}}
 \prove{$S$ is star-closed.}
\end{cgss}
\begin{proof}
 \pfsketch\ We will orient the two lines so that their inner product is $-1$ and test their
 sum against every root. This is the angle argument in the paper, including coincident lines.
 \step{23sum}{\assume{$\Zr,\Zs\in\Phi$ and $(\Zr,\Zs)=-1$.}
  \prove{$\Zr+\Zs\in\Phi$.}}
 \begin{proof}
  \step{23norm}{$(\Zr+\Zs,\Zr+\Zs)=2$.}
  \begin{proof}
   \pf\ By bilinearity and \asm{1}, the squared norm is $2-1-1+2=2$.~\qed
  \end{proof}
  \step{23test}{\assume{$\Zt\in\Phi$.}
   \prove{$(\Zr+\Zs,\Zt)\in\{-2,-1,0,1,2\}$.}}
  \begin{proof}
   \step{23same}{\case{$\Zt\in\{\Zr,-\Zr,\Zs,-\Zs\}$.}}
   \begin{proof}
    \pf\ The four values are $1,-1,1,-1$, respectively, by bilinearity and \asm{1}.~\qed
   \end{proof}
   \step{23other}{\case{$\Zt\notin\{\Zr,-\Zr,\Zs,-\Zs\}$.}}
   \begin{proof}
    \step{23small}{$(\Zr,\Zt),(\Zs,\Zt)\in\{-1,0,1\}$.}
    \begin{proof}
     \pf\ The line-system condition and \eqref{eq:endpoints} exclude $\pm2$.~\qed
    \end{proof}
    \qedstep
    \begin{proof}
     \pf\ The sum of two integers between $-1$ and 1 is an integer between $-2$ and 2;
     apply bilinearity to \stepref{23small}.~\qed
    \end{proof}
   \end{proof}
   \qedstep
   \begin{proof}\pf\ The cases \stepref{23same} and \stepref{23other} are exhaustive.~\qed\end{proof}
  \end{proof}
  \step{23span}{$\Zr+\Zs\in\operatorname{span}(\Phi)$.}
  \begin{proof}
   \pf\ Both $\Zr$ and $\Zs$ lie in $\Phi$ by \asm{1}, and a span is closed under
   addition.~\qed
  \end{proof}
  \qedstep
  \begin{proof}\pf\ Apply Definition 2.2 to \stepref{23span}, \stepref{23norm} and
   \stepref{23test}.~\qed\end{proof}
 \end{proof}
 \qedstep
 \begin{proof}\pf\ Step \stepref{23sum} is Definition 2.1.~\qed\end{proof}
\end{proof}

\begin{cgss}{Hypothesis}{2.4}\label{hyp:2.4}
 \lean{LineSystem.Hyp24}\uses{def:line-system,def:indecomposable,def:2.1}
 \assume{\begin{pfenum}
  \item $V$ is a finite-dimensional real inner product space and $S$ is a line system in $V$.
  \item $S$ is indecomposable and star-closed.
  \item $\Za,\Zb\in\Phi$ and $(\Za,\Zb)=-1$.
 \end{pfenum}}
 \pflet{$\Zc\triangleq-\Za-\Zb$, and}
 \begin{align*}
 \Ga&\triangleq\{\Zx\in\Phi:(\Zx,\Za)=0,\ (\Zx,\Zb)=1,\ (\Zx,\Zc)=-1\},\\
 \Gb&\triangleq\{\Zx\in\Phi:(\Zx,\Zb)=0,\ (\Zx,\Zc)=1,\ (\Zx,\Za)=-1\},\\
 \Gc&\triangleq\{\Zx\in\Phi:(\Zx,\Zc)=0,\ (\Zx,\Za)=1,\ (\Zx,\Zb)=-1\},\\
 \Delta&\triangleq\{\Zt\in\Phi:(\Zt,\Za)=(\Zt,\Zb)=(\Zt,\Zc)=0\},\\
 \Sigma&\triangleq\{\Za,\Zb,\Zc\}\cup\Ga\cup\Gb\cup\Gc\cup\Delta.
 \end{align*}
 We will write $X=\opk{Orth}(\Ga)$ for the \emph{orthogonality graph} on $\Ga$: distinct vertices
 $x,y$ are adjacent exactly when $(\Zx,\Zy)=0$. We will write a vertex as $x$ and
 its representing vector as $\Zx$.
 \label{def:gamma-graph}\lean{LineSystem.Hyp24.orthGraph}\uses{hyp:2.4}
 The represented graph of a vector family uses inner product 1; its
 \emph{obtuse graph}, denoted $\opk{Obt}$, uses inner product $-1$.
 We use the orthogonality convention throughout Section 2.
 The \emph{empty graph} $\comp{K_m}$ has $m$ vertices and no edges, with $m=0$ allowed.
\end{cgss}

\noindent\repair\ We require the pair $\Za,\Zb$ explicitly (Implicit hypothesis 4 of
\path{cgss/errata.md}). A one-line system has no such pair;
we will treat it separately in Lemma 2.12 and Theorem 2.13.
\repair\ Also, $\Delta=-\Delta$ here (Error 6), so
$\Sigma$ need not have just one vector on each line. Lemma 2.5 asserts a signed cover.

\begin{namedresult}{Lemma}{(star identities)}\label{lem:star-identities}
 \lean{LineSystem.Hyp24.c_mem}\lean{LineSystem.Hyp24.inner_b_c}
 \lean{LineSystem.Hyp24.inner_c_a}\lean{LineSystem.Hyp24.inner_c_b}
 \uses{hyp:2.4,def:2.1,facts:inner}
 \assume{Hypothesis 2.4 holds in the finite-dimensional real inner product space $V$.}
 \prove{$\Zc\in\Phi$ and
  $(\Za,\Zb)=(\Zb,\Zc)=(\Zc,\Za)=(\Zc,\Zb)=-1$.}
\end{namedresult}
\begin{proof}
 \pfsketch\ \new\ We will state the star identities separately so that later results
 can cite them without citing a step inside the proof of Lemma 2.5.
 \step{25star}{$\Zc\in\Phi$ and all three inner products between $\Za,\Zb,\Zc$ are $-1$.}
 \begin{proof}
  \step{25mem}{$\Za+\Zb\in\Phi$, so $\Zc\in\Phi$.}
  \begin{proof}\pf\ Use star-closure, $(\Za,\Zb)=-1$, and closure under negation.~\qed\end{proof}
  \qedstep
  \begin{proof}
   \pf\ By bilinearity, $(\Za,\Zc)=-2+1=-1$ and $(\Zb,\Zc)=1-2=-1$;
   the remaining products follow from symmetry and the hypothesis $(\Za,\Zb)=-1$.~\qed
  \end{proof}
 \end{proof}
 \qedstep
 \begin{proof}\pf\ Step \stepref{25star} and symmetry give the stated identities.~\qed\end{proof}
\end{proof}

\begin{cgss}{Lemma}{2.5}\label{lem:2.5}
 \lean{LineSystem.Hyp24.Φ_eq_sigmaSet_union_neg}
 \uses{hyp:2.4,lem:star-identities,facts:inner}
 \assume{Hypothesis 2.4 holds in the finite-dimensional real inner product space $V$.}
 \prove{$\Phi=\Sigma\cup(-\Sigma)$; equivalently, $S$ consists exactly of the lines of $\Sigma$.}
\end{cgss}
\begin{proof}
 \pfsketch\ We will exclude equal nonzero inner products with $\Za,\Zb$, then choose the
 sign in each of the remaining seven cases.
 \step{25equal}{\assume{$\Zu\in\Phi$, $\epsilon\in\{-1,1\}$, and
   $(\Zu,\Za)=(\Zu,\Zb)=\epsilon$.}
  \prove{$\Zu=-\epsilon\Zc$.}}
 \begin{proof}
  \step{25zero}{$(\Za+\Zb-\epsilon\Zu,\Za+\Zb-\epsilon\Zu)=0$.}
  \begin{proof}
   \pf\ Bilinearity gives $2+2\epsilon^2-2\epsilon(2\epsilon)=0$, since
   $\epsilon^2=1$ and $(\Za+\Zb,\Za+\Zb)=2$.~\qed
  \end{proof}
  \qedstep
  \begin{proof}
   \pf\ Zero squared norm gives $\Za+\Zb=\epsilon\Zu$; multiply by $\epsilon$ and use
   $\Zc=-\Za-\Zb$ and $\epsilon^2=1$.~\qed
  \end{proof}
 \end{proof}
 \step{25switch}{\assume{$\Zv\in\Phi$.}
  \prove{$\Zv\in\Sigma$ or $-\Zv\in\Sigma$.}}
 \begin{proof}
  \step{25onstar}{\case{$\Zv\in\{\pm\Za,\pm\Zb,\pm\Zc\}$.}}
  \begin{proof}\pf\ Choose the displayed sign of a member of $\{\Za,\Zb,\Zc\}\subseteq\Sigma$.~\qed\end{proof}
  \step{25offstar}{\case{$\Zv\notin\{\pm\Za,\pm\Zb,\pm\Zc\}$.}}
  \begin{proof}
   \step{25poss}{$((\Zv,\Za),(\Zv,\Zb))$ is one of
    $(0,0),(0,1),(0,-1),(1,0),(-1,0),(1,-1),(-1,1)$.}
   \begin{proof}
    \pf\ By \eqref{eq:endpoints}, both entries are in $\{-1,0,1\}$.
    Step \stepref{25equal} excludes the two pairs $(1,1)$ and $(-1,-1)$.~\qed
   \end{proof}
   \step{25table}{We choose representatives by the following table.}
   \begin{proof}
    \pf\ In each row $(\Zv,\Zc)=-(\Zv,\Za)-(\Zv,\Zb)$ by bilinearity;
    substitution gives the defining three inner products of the indicated set.
    \[
    \begin{array}{c|rrrrrrr}
    ((\Zv,\Za),(\Zv,\Zb))&(0,0)&(0,1)&(0,-1)&(1,0)&(-1,0)&(1,-1)&(-1,1)\\\hline
    \text{representative}&\Zv&\Zv&-\Zv&-\Zv&\Zv&\Zv&-\Zv\\
    \text{set}&\Delta&\Ga&\Ga&\Gb&\Gb&\Gc&\Gc
    \end{array}
    \]
    Each representative belongs to $\Phi$ by closure under negation.~\qed
   \end{proof}
   \qedstep
   \begin{proof}\pf\ Apply \stepref{25table} to the exhaustive list \stepref{25poss}.~\qed\end{proof}
  \end{proof}
  \qedstep
  \begin{proof}\pf\ The cases \stepref{25onstar} and \stepref{25offstar} are exhaustive.~\qed\end{proof}
 \end{proof}
 \qedstep
 \begin{proof}
  \pf\ Step \stepref{25switch} gives $\Phi\subseteq\Sigma\cup(-\Sigma)$.
  Conversely, the star identities and the definitions give $\Sigma\subseteq\Phi$,
  and $\Phi=-\Phi$ gives $-\Sigma\subseteq\Phi$.~\qed
 \end{proof}
\end{proof}

\begin{cgss}{Lemma}{2.6}\label{lem:2.6}
 \lean{LineSystem.Hyp24.ΓaEquivΓb}
 \lean{LineSystem.Hyp24.ΓbEquivΓc}\lean{LineSystem.Hyp24.ΓcEquivΓa}\lean{LineSystem.Hyp24.inner_eq_zero_or_one}
 \lean{LineSystem.Hyp24.add_c_mem_Γb}\lean{LineSystem.Hyp24.sub_b_mem_Γc}
 \lean{LineSystem.Hyp24.inner_ΓaEquivΓb}\lean{LineSystem.Hyp24.inner_ΓbEquivΓc}
 \lean{LineSystem.Hyp24.inner_ΓcEquivΓa}
 \lean{LineSystem.Hyp24.inner_eq_zero_or_one_of_mem_Γb}
 \lean{LineSystem.Hyp24.inner_eq_zero_or_one_of_mem_Γc}
 \uses{hyp:2.4,lem:star-identities,facts:inner}
 \assume{Hypothesis 2.4 holds in the finite-dimensional real inner product space $V$.}
 \prove{\begin{pfenum}
  \item The maps $\varphi_c(\Zx)=\Zx+\Zc$, $\varphi_a(\Zy)=\Zy+\Za$, and
   $\varphi_b(\Zz)=\Zz+\Zb$ give inner-product-preserving bijections
   $\Ga\to\Gb\to\Gc\to\Ga$.
  \item $\Gb=\Zc+\Ga$ and $\Gc=-\Zb+\Ga$.
  \item For each $D\in\{\Ga,\Gb,\Gc\}$ and distinct $\Zx,\Zy\in D$,
   $(\Zx,\Zy)\in\{0,1\}$. Diagonal inner products are 2.
 \end{pfenum}}
\end{cgss}
\begin{proof}
 \pfsketch\ We will prove the translation identity for all three cyclic triples at once.
 Their composite is the identity, so each map is bijective. A negative inner product would
 give opposite vectors in two incompatible parts of the partition.
 \step{26cycle}{\assume{$(\Zr,\Zs,\Zt)$ is one of
  $(\Za,\Zb,\Zc),(\Zb,\Zc,\Za),(\Zc,\Za,\Zb)$, and $\Zx\in\Gamma_r$.}
  \prove{$\Zx+\Zt\in\Gamma_s$.}}
 \begin{proof}
  \step{26root}{$\Zx+\Zt\in\Phi$.}
  \begin{proof}
   \pf\ In each of the three listed triples, $(\Zx,\Zt)=-1$ by the definition of
   $\Gamma_r$; apply star-closure and the star identities.~\qed
  \end{proof}
  \step{26values}{$(\Zx+\Zt,\Zr)=-1$, $(\Zx+\Zt,\Zs)=0$, and
   $(\Zx+\Zt,\Zt)=1$.}
  \begin{proof}
   \pf\ By bilinearity, the star identities and the defining products, these values are respectively
   $0-1$, $1-1$, and $-1+2$.~\qed
  \end{proof}
  \qedstep
  \begin{proof}\pf\ Steps \stepref{26root} and \stepref{26values} are the definition of $\Gamma_s$.~\qed\end{proof}
 \end{proof}
 \step{26bijections}{The three translations are bijections and preserve inner products.}
 \begin{proof}
  \step{26inverse}{Each translation has the next two translations as a two-sided inverse.}
  \begin{proof}
   \pf\ By \stepref{26cycle} all composites have the stated domains, and every cyclic
   composite sends $\Zx$ to $\Zx+\Za+\Zb+\Zc=\Zx$.~\qed
  \end{proof}
  \step{26gram}{\assume{$(\Zr,\Zs,\Zt)$ is a listed cyclic triple and
    $\Zx,\Zy\in\Gamma_r$.}
   \prove{$(\Zx+\Zt,\Zy+\Zt)=(\Zx,\Zy)$.}}
  \begin{proof}
   \pf\ By bilinearity the left side is $(\Zx,\Zy)-1-1+2=(\Zx,\Zy)$,
   using the definition of $\Gamma_r$.~\qed
  \end{proof}
  \qedstep
  \begin{proof}\pf\ Combine the two-sided inverses in \stepref{26inverse} with \stepref{26gram}.~\qed\end{proof}
 \end{proof}
 \step{26positive}{\assume{$(\Zr,\Zs,\Zt)$ is a listed cyclic triple and
    $\Zx,\Zy\in\Gamma_r$ are distinct.}
  \prove{$(\Zx,\Zy)\in\{0,1\}$.}}
 \begin{proof}
  \step{26endpoints}{$(\Zx,\Zy)\ne2,-2$.}
  \begin{proof}
   \pf\ The value 2 contradicts distinctness by \eqref{eq:endpoints}.
   The value $-2$ gives $\Zy=-\Zx$, contradicting
   $(\Zx,\Zs)=(\Zy,\Zs)=1$.~\qed
  \end{proof}
  \step{26minus}{$(\Zx,\Zy)\ne-1$.}
  \begin{proof}
   \step{26opposite}{\assume{$(\Zx,\Zy)=-1$.}
    \prove{$\Zx=-(\Zy+\Zt)$.}}
   \begin{proof}
    \pf\ Step \stepref{26cycle} puts $\Zy+\Zt$ in $\Phi$, and bilinearity gives
    $(\Zx,\Zy+\Zt)=-1-1=-2$; apply \eqref{eq:endpoints}.~\qed
   \end{proof}
   \qedstep
   \begin{proof}
    \pf\ The equality in \stepref{26opposite} gives
    $0=(\Zx,\Zr)=-(\Zy+\Zt,\Zr)=1$, by the definition of $\Gamma_r$
    and $(\Zt,\Zr)=-1$ from the star identities, a contradiction.~\qed
   \end{proof}
  \end{proof}
  \qedstep
  \begin{proof}\pf\ The line-system condition, \stepref{26endpoints}, and \stepref{26minus} leave 0 and 1.~\qed\end{proof}
 \end{proof}
 \qedstep
 \begin{proof}
  \pf\ Step \stepref{26bijections} proves part 1. Its first map gives
  $\Gb=\Zc+\Ga$, and its first two maps give
  $\Gc=(\Za+\Zc)+\Ga=-\Zb+\Ga$. Step \stepref{26positive} and the norm condition prove part 3.~\qed
 \end{proof}
\end{proof}

\begin{cgss}{Lemma}{2.7}\label{lem:2.7}
 \lean{LineSystem.Hyp24.Δ_eq_sub}\lean{LineSystem.Hyp24.Φ_eq_of_Γa_eq}
 \lean{LineSystem.Hyp24.Φ_eq_recovery_union_neg}\lean{LineSystem.Hyp24.sigmaSet_eq_recoverySet}
 \uses{hyp:2.4,lem:2.5,lem:2.6,def:indecomposable,facts:inner}
 \assume{Hypothesis 2.4 holds in the finite-dimensional real inner product space $V$.}
 \prove{\[
  \Delta=\{\Zx-\Zy:\Zx,\Zy\in\Ga,\ (\Zx,\Zy)=1\}.
 \]
 Put
 \[
 \cR=\{\Za,\Zb,\Zc\}\cup\Ga\cup(\Zc+\Ga)\cup(-\Zb+\Ga)
       \cup\{\Zx-\Zy:\Zx,\Zy\in\Ga,\ (\Zx,\Zy)=1\}.
 \]
 Then $\Sigma=\cR$ and $\Phi=\cR\cup(-\cR)$.
 In particular, any two instances of Hypothesis 2.4 in $V$ with the same
 $\Za,\Zb$ and $\Ga$ have the same signed root set $\Phi$.}
\end{cgss}
\begin{proof}
 \pfsketch\ We will recover $\Gb,\Gc$ by translation and $\Delta$ by differences.
 A root of $\Delta$ orthogonal to all of $\Ga$ would belong to an orthogonal part of $\Phi$;
 we will give that partition explicitly.
 \step{27forward}{\assume{$\Zx,\Zy\in\Ga$ and $(\Zx,\Zy)=1$.}
  \prove{$\Zx-\Zy\in\Delta$.}}
 \begin{proof}
  \step{27root}{$\Zx-\Zy\in\Phi$.}
  \begin{proof}\pf\ Apply star-closure to $\Zx,-\Zy\in\Phi$, whose inner product is $-1$.~\qed\end{proof}
  \qedstep
  \begin{proof}
   \pf\ The products of $\Zx-\Zy$ with $\Za,\Zb,\Zc$ are
   $0-0,1-1,-1-(-1)$, respectively, by bilinearity. With \stepref{27root}, this defines $\Delta$.~\qed
  \end{proof}
 \end{proof}
 \step{27nonzero}{\assume{$\Zt\in\Delta$ and $(\Zt,\Zx)\ne0$ for some $\Zx\in\Ga$.}
  \prove{$\Zt=\Zu-\Zv$ for some $\Zu,\Zv\in\Ga$ with $(\Zu,\Zv)=1$.}}
 \begin{proof}
  \pflet{Choose $\Zx\in\Ga$ with $(\Zt,\Zx)\ne0$.}
  \step{27sign}{$(\Zt,\Zx)\in\{-1,1\}$.}
  \begin{proof}
   \pf\ If the product were $\pm2$, \eqref{eq:endpoints} would give
   $\Zt=\pm\Zx$, contradicting $(\Zt,\Zb)=0$ and $(\Zx,\Zb)=1$.
   The line-system condition and the nonzero assumption leave $\pm1$.~\qed
  \end{proof}
  \step{27negative}{\case{$(\Zt,\Zx)=-1$.}}
  \begin{proof}
   \step{27add}{$\Zx+\Zt\in\Ga$.}
   \begin{proof}
    \pf\ Star-closure gives membership in $\Phi$; adding $\Zt$, which is orthogonal to
    $\Za,\Zb,\Zc$, preserves the three products defining $\Ga$.~\qed
   \end{proof}
   \qedstep
   \begin{proof}
    \pf\ Choose $\Zu=\Zx+\Zt$, $\Zv=\Zx$. Their difference is $\Zt$, and
    $(\Zu,\Zv)=2-1=1$ by bilinearity and \asm{2}.~\qed
   \end{proof}
  \end{proof}
  \step{27positive}{\case{$(\Zt,\Zx)=1$.}}
  \begin{proof}
   \step{27sub}{$\Zx-\Zt\in\Ga$.}
   \begin{proof}
    \pf\ Star-closure applies to $\Zx,-\Zt$; subtracting $\Zt$ preserves the
    three products defining $\Ga$.~\qed
   \end{proof}
   \qedstep
   \begin{proof}
    \pf\ Choose $\Zu=\Zx$, $\Zv=\Zx-\Zt$. Their difference is $\Zt$, and
    $(\Zu,\Zv)=2-1=1$ by bilinearity and \asm{2}.~\qed
   \end{proof}
  \end{proof}
  \qedstep
  \begin{proof}\pf\ Step \stepref{27sign} makes \stepref{27negative} and \stepref{27positive} exhaustive.~\qed\end{proof}
 \end{proof}
 \step{27partition}{No root of $\Delta$ is orthogonal to all of $\Ga$.}
 \begin{proof}
  \pflet{$P\triangleq\{\Zt\in\Delta:\forall\Zx\in\Ga,\ (\Zt,\Zx)=0\}$ and
   $Q\triangleq\Phi\diff P$.}
  \step{27signstable}{$P=-P$ and $\Za\in Q$.}
  \begin{proof}
   \pf\ The defining products of $P$ are all zero, and thus invariant under negation.
   Also $\Za\in\Phi$ and $(\Za,\Za)=2$, so $\Za\notin\Delta$.~\qed
  \end{proof}
  \step{27sigma}{\assume{$\Zt\in P$ and $\Zu\in\Sigma\diff P$.}
   \prove{$(\Zt,\Zu)=0$.}}
  \begin{proof}
   \step{27core}{\case{$\Zu\in\{\Za,\Zb,\Zc\}\cup\Ga\cup\Gb\cup\Gc$.}}
   \begin{proof}
    \pf\ The definition of $P$ gives orthogonality to $\Za,\Zb,\Zc,\Ga$.
    Lemma 2.6 gives $\Gb=\Zc+\Ga$ and $\Gc=-\Zb+\Ga$; bilinearity gives zero
    on each of these translates.~\qed
   \end{proof}
   \step{27delta}{\case{$\Zu\in\Delta\diff P$.}}
   \begin{proof}
    \step{27witness}{$\Zu=\Zx-\Zy$ for some $\Zx,\Zy\in\Ga$.}
    \begin{proof}
     \pf\ Membership in $\Delta\diff P$ gives a nonzero product with a member of $\Ga$;
     apply \stepref{27nonzero}.~\qed
    \end{proof}
    \qedstep
    \begin{proof}\pf\ By the definition of $P$, $(\Zt,\Zx-\Zy)=0-0=0$.~\qed\end{proof}
   \end{proof}
   \qedstep
   \begin{proof}\pf\ The definition of $\Sigma$ makes \stepref{27core} and \stepref{27delta} exhaustive.~\qed\end{proof}
  \end{proof}
  \step{27orthogonal}{$(\Zt,\Zu)=0$ for all $\Zt\in P$ and $\Zu\in Q$.}
  \begin{proof}
   \pf\ Lemma 2.5 gives $\Zu=\epsilon\Zv$ with $\epsilon\in\{-1,1\}$ and
   $\Zv\in\Sigma$. By \stepref{27signstable}, $\Zv\notin P$.
   Step \stepref{27sigma} and bilinearity give $(\Zt,\Zu)=\epsilon\cdot0=0$.~\qed
  \end{proof}
  \qedstep
  \begin{proof}
   \pf\ We have $P\cup Q=\Phi$. Indecomposability and \stepref{27orthogonal} imply
   $P=\varnothing$ or $Q=\varnothing$; \stepref{27signstable} excludes the second alternative.~\qed
  \end{proof}
 \end{proof}
 \qedstep
 \begin{proof}
  \pf\ Steps \stepref{27forward}, \stepref{27nonzero}, and \stepref{27partition} give both
  inclusions in the formula for $\Delta$. Substitute this equality and Lemma 2.6
  into the definition of $\Sigma$ to obtain $\Sigma=\cR$; Lemma 2.5 gives
  $\Phi=\cR\cup(-\cR)$. Equal $\Za,\Zb,\Ga$ give equal $\Zc=-\Za-\Zb$
  and equal $\cR$, so they give equal root sets.~\qed
 \end{proof}
\end{proof}

\begin{cgss}{Lemma}{2.8}\label{lem:2.8}
  \lean{LineSystem.Hyp24.triangle}\uses{hyp:2.4,lem:2.6,lem:star-identities,facts:inner}
  \assume{\begin{pfenum}
      \item Hypothesis 2.4 holds in the finite-dimensional real inner product space $V$.
      \item $\Zx,\Zy\in\Ga$ and $(\Zx,\Zy)=0$.
    \end{pfenum}}
  \prove{\begin{pfenum}
      \item $\Zz\triangleq\Zb-\Zc-\Zx-\Zy$ lies in $\Ga$.
      \item Each $\Zw\in\Ga\diff\{\Zx,\Zy,\Zz\}$ is orthogonal to exactly one of $\Zx$, $\Zy$,
        $\Zz$.
    \end{pfenum}}
\end{cgss}

\begin{proof}
  \pfsketch\ The vectors $\Zx+\Zc$ and $\Zy-\Zb$ lie in $\Phi$ and have inner product $-1$, so
  star-closure puts their sum, which is $-\Zz$, in $\Phi$. Computing inner products places $\Zz$
  in $\Ga$. For part 2, $(\Zw,\Zx)+(\Zw,\Zy)+(\Zw,\Zz)=2$, and each term is 0 or 1.

  \step{xc}{$\Zx+\Zc\in\Gb$ and $\Zy-\Zb\in\Gc$.}
  \begin{proof}
    \pf\ By \asm[2]{0} and Lemma 2.6.~\qed
  \end{proof}

  \step{ip}{$(\Zx+\Zc,\Zy-\Zb)=-1$.}
  \begin{proof}
    \pf\ We use $(\Zb,\Zc)=-1$ from the star identities.
    \begin{align*}
      (\Zx+\Zc,\Zy-\Zb)&=(\Zx,\Zy)-(\Zx,\Zb)+(\Zc,\Zy)-(\Zc,\Zb)
        &&\text{[bilinearity]}\\
      &=0-1-1+1
        &&\text{[\asm[2]{0}, the definition of $\Ga$, and $(\Zb,\Zc)=-1$]}\\
      &=-1.
    \end{align*}
    \qed
  \end{proof}

  \step{zphi}{$\Zz\in\Phi$.}
  \begin{proof}
    \pf\ By \stepref{xc}, \stepref{ip} and star-closure, $(\Zx+\Zc)+(\Zy-\Zb)=-\Zz$ lies in
    $\Phi$; and $\Phi=-\Phi$.~\qed
  \end{proof}

  \step{zga}{$\Zz\in\Ga$.}
  \begin{proof}
    \pf\ We have $(\Zz,\Za)=0$, $(\Zz,\Zb)=1$ and $(\Zz,\Zc)=-1$, by bilinearity, \asm[2]{0} and
    the definition of $\Ga$; together with \stepref{zphi}, this is the definition of
    $\Ga$.~\qed
  \end{proof}

  \step{w}{\assume{$\Zw\in\Ga\diff\{\Zx,\Zy,\Zz\}$.}
    \prove{$\Zw$ is orthogonal to exactly one of $\Zx$, $\Zy$, $\Zz$.}}
  \begin{proof}
    \step{sum}{$(\Zw,\Zx)+(\Zw,\Zy)+(\Zw,\Zz)=2$.}
    \begin{proof}
      \pf\
      \begin{align*}
        (\Zw,\Zx)+(\Zw,\Zy)+(\Zw,\Zz)&=(\Zw,\Zb-\Zc)&&\text{[the definition of $\Zz$]}\\
        &=1+1&&\text{[\asm{1} and the definition of $\Ga$]}
      \end{align*}
      \qed
    \end{proof}

    \step{zo}{Each of $(\Zw,\Zx)$, $(\Zw,\Zy)$, $(\Zw,\Zz)$ is 0 or 1.}
    \begin{proof}
      \pf\ By \asm{1}, \asm[2]{0}, \stepref{zga} and Lemma 2.6.~\qed
    \end{proof}

    \qedstep
    \begin{proof}
      \pf\ By \stepref{sum} and \stepref{zo}, three numbers in $\{0,1\}$ have sum 2, so exactly
      one of them is 0.~\qed
    \end{proof}
  \end{proof}

  \qedstep
  \begin{proof}
    \pf\ By \stepref{zga} and \stepref{w}.~\qed
  \end{proof}
\end{proof}

\begin{namedresult}{Definition}{(the triangle property)}\label{def:triangle-property}
 \lean{SimpleGraph.HasTriangleProperty}\uses{}
 \assume{$X$ is a simple graph on a set $U$.}
 We say that $X$ has the \emph{triangle property} if every edge $xy$ has exactly one
 common neighbour $z$, and every vertex outside $\{x,y,z\}$ is adjacent to exactly one
 of $x,y,z$. For adjacent $x,y$ we will write $f(x,y)$ for this unique common neighbour;
 the graph will be clear from the argument.
 \new\ We name this predicate once so that we can use the graph argument without a vector
 representation.
\end{namedresult}

\begin{namedresult}{Corollary}{of Lemma 2.8}\label{cor:triangle-property}
 \lean{LineSystem.Hyp24.orthGraph_hasTriangleProperty}\uses{hyp:2.4,lem:2.8,def:triangle-property,facts:inner}
 \assume{Hypothesis 2.4 holds in the finite-dimensional real inner product space $V$, and $X$ is its orthogonality graph on $\Ga$.}
 \prove{$X$ has the triangle property.}
\end{namedresult}
\begin{proof}
 \pfsketch\ \new\ We will record the pairwise orthogonality calculation needed to pass
 from Lemma 2.8 to graph terminology; it does not change that lemma.
 \step{28triangle}{\assume{$x,y$ are adjacent vertices of $X$.}
  \prove{$\Zz\triangleq\Zb-\Zc-\Zx-\Zy$ represents a distinct common neighbour $z$.}}
 \begin{proof}
  \step{28products}{$(\Zx,\Zz)=(\Zy,\Zz)=0$.}
  \begin{proof}
   \pf\ Bilinearity gives $(\Zx,\Zz)=1-(-1)-2-0=0$ and
   $(\Zy,\Zz)=1-(-1)-0-2=0$, using adjacency and the definition of $\Ga$.~\qed
  \end{proof}
  \qedstep
  \begin{proof}
   \pf\ Lemma 2.8 puts $\Zz$ in $\Ga$. Step \stepref{28products} and squared norm 2
   give $z\ne x,y$ and $z\sim x,y$.~\qed
  \end{proof}
 \end{proof}
 \step{28unique}{\assume{$w$ is a common neighbour of adjacent $x,y$.}
  \prove{$\Zw=\Zb-\Zc-\Zx-\Zy$.}}
 \begin{proof}
  \pf\ If $w$ differed from the vertex in \stepref{28triangle}, it would lie outside
  its triangle and be adjacent to two of its vertices, contrary to Lemma 2.8.~\qed
 \end{proof}
 \qedstep
 \begin{proof}
  \pf\ Steps \stepref{28triangle} and \stepref{28unique} give the unique common neighbour;
  Lemma 2.8 gives the assertion about every outside vertex.~\qed
 \end{proof}
\end{proof}

\begin{cgss}{Lemma}{2.9}\label{lem:2.9}
 \lean{LineSystem.Hyp24.orthGraph_iso_friendship_of_forall_adj}\uses{hyp:2.4,def:triangle-property,cor:triangle-property}
 \assume{\begin{pfenum}
  \item Hypothesis 2.4 holds in the finite-dimensional real inner product space $V$.
  \item $X$ is its orthogonality graph and $x\in V(X)$ is adjacent to every other vertex.
 \end{pfenum}}
 \prove{There is an integer $m\geq0$ and a partition of $V(X)\diff\{x\}$ into $m$ pairs,
  each of which makes a triangle with $x$. These are all the edges of $X$; thus $X$ is the friendship graph with $m$ triangles.}
\end{cgss}
\begin{proof}
 \pfsketch\ The unique triangle on each edge through $x$ pairs the other vertices.
 We will use the outside-vertex condition to exclude edges between different pairs.
 \step{29mate}{For $y\ne x$, let $f(x,y)$ be the unique common neighbour of $x,y$.
  Then $f(x,y)\ne y,x$ and $f(x,f(x,y))=y$.}
 \begin{proof}
  \step{29exists}{$f(x,y)$ exists and differs from $x,y$.}
  \begin{proof}\pf\ Apply Corollary~\ref{cor:triangle-property} to the edge $xy$.~\qed\end{proof}
  \qedstep
  \begin{proof}
   \pf\ The vertex $y$ is a common neighbour of $x,f(x,y)$, so uniqueness in
   Corollary~\ref{cor:triangle-property} gives $f(x,f(x,y))=y$.~\qed
  \end{proof}
 \end{proof}
 \step{29pairs}{The sets $\{y,f(x,y)\}$ partition $V(X)\diff\{x\}$ into pairs.}
 \begin{proof}
  \pf\ Step \stepref{29mate} gives a fixed-point-free involution: its two-element
  orbits cover the domain, and intersecting orbits coincide by applying the involution.~\qed
 \end{proof}
 \step{29edges}{There is no edge between two different pairs.}
 \begin{proof}
  \pf\ A vertex in a different pair lies outside $\{x,y,f(x,y)\}$ and is adjacent to $x$.
  Corollary~\ref{cor:triangle-property} excludes adjacency to either $y$ or $f(x,y)$.~\qed
 \end{proof}
 \qedstep
 \begin{proof}
  \pf\ The finite partition in \stepref{29pairs} has some size $m\geq0$.
  Step \stepref{29mate} supplies the triangle edges, and \stepref{29edges} excludes all others.
  When $V(X)=\{x\}$ the partition is empty and $m=0$.~\qed
 \end{proof}
\end{proof}

\begin{namedresult}{Lemma}{(triangle-property normal forms)}\label{lem:triangle-normal-form}
 \lean{SimpleGraph.HasTriangleProperty.normal_form_of_not_forall_adj}
 \uses{def:triangle-property}
 \assume{\begin{pfenum}
  \item $X$ is a finite simple graph with the triangle property.
  \item $X$ has an edge.
  \item No vertex of $X$ is adjacent to every other vertex.
 \end{pfenum}}
 \prove{For every pair of distinct nonadjacent vertices $u,v$, put
 \[
 \begin{array}{ll}
 K=\Gamma_1(u)\cap\Gamma_1(v),&L=\Gamma_1(u)\diff\Gamma_1(v),\\
 M=\Gamma_1(v)\diff\Gamma_1(u),&
 N=V(X)\diff\bigl(\{u,v\}\cup K\cup L\cup M\bigr),
 \end{array}
 \]
 and put $k=|K|$. Then the following hold.
 \begin{pfenum}
  \item $k\in\{2,3,5\}$, $X$ is regular of degree $2k$, $|V(X)|=6k-3$, and
   \repair\ $X$ is strongly regular with parameters $(6k-3,2k,1,k)$
   (Error 12 of \path{cgss/errata.md}).
  \item The sets $\{u\},\{v\},K,L,M,N$ partition $V(X)$, with $|L|=|M|=k$ and
   $|N|=3k-5$. The labels $x'=f(u,x)$ and $x''=f(v,x)$ give bijections
   $K\to L$ and $K\to M$.
  \item $K,L,M$ are independent. For $x,y\in K$, $x\sim y'$ and $x\sim y''$
   hold exactly when $x=y$, and $x'\sim y''$ holds exactly when $x\ne y$.
  \item For $w\in N$, put $A(w)=K\diff\Gamma_1(w)$ and
   $d_K(w)=|K\cap\Gamma_1(w)|$. Then $|A(w)|\in\{0,2\}$,
   $d_K(w)\in\{k-2,k\}$, and $|N\cap\Gamma_1(w)|=d_K(w)$.
   For $x\in K$, $w\sim x'$ and $w\sim x''$ hold exactly when $x\in A(w)$.
  \item The graph on $N$ has $k(k-2)$ edges, each with its third vertex in $K$.
   Moreover,
   \[
    \sum_{w\in N}d_K(w)=k(2k-4),\qquad
    2|\{w\in N:d_K(w)=k\}|=(k-2)(5-k).
   \]
  \item If $k=2$, $N$ consists of one vertex $w$, with $A(w)=K$, and has no edges.
  \item If $k=3$, the graph on $N$ is a claw with centre $w_0$, where $A(w_0)=\varnothing$.
   Its leaves are $w_x$ for $x\in K$, with $A(w_x)=K\diff\{x\}$.
  \item If $k=5$, the map $w\mapsto A(w)$ is a bijection from $N$ to the two-element
   subsets of $K$, and adjacency in $N$ is disjointness of these labels.
 \end{pfenum}
 These rules and the definitions of $K,L,M,N$ specify every adjacency.}
\end{namedresult}
\begin{proof}
 \pfsketch\ \new\ We separate the combinatorial argument from its application to
 line systems. We will keep the paper's partition and prove regularity for arbitrary
 pairs. The invariant $d_K$ gives both the bound on $k$ and the three normal forms.
 All counting identities involving subtraction are read in $\ints$.

 \step{210triangles}{Every edge $xy$ has a unique common neighbour $f(x,y)$.
  Every vertex has positive degree and has a distinct non-neighbour.}
 \begin{proof}
  \step{210third}{The triangle on $xy$ is $\{x,y,f(x,y)\}$ and every outside vertex
   sees exactly one vertex of it.}
  \begin{proof}\pf\ Apply the triangle property in \asm[1]{0}.~\qed\end{proof}
  \step{210positive}{Every vertex has positive degree.}
  \begin{proof}
   \pf\ A triangle exists since $X$ has an edge. Its vertices have degree at least 2,
   and every other vertex has a neighbour in it by \stepref{210third}.~\qed
  \end{proof}
  \step{210nouniversal}{Every vertex has a distinct non-neighbour.}
  \begin{proof}
   \pf\ This is precisely \asm[3]{0}, the absence of a universal vertex.~\qed
  \end{proof}
  \qedstep
  \begin{proof}\pf\ Combine \stepref{210third}, \stepref{210positive}, and \stepref{210nouniversal}.~\qed\end{proof}
 \end{proof}

 \step{210local}{For every pair of distinct nonadjacent vertices $u,v$, define, excluding $u,v$,
  \[
  \begin{array}{ll}
  K=\Gamma_1(u)\cap\Gamma_1(v),&
  L=\Gamma_1(u)\diff\Gamma_1(v),\\
  M=\Gamma_1(v)\diff\Gamma_1(u),&
  N=V(X)\diff\bigl(\{u,v\}\cup K\cup L\cup M\bigr).
  \end{array}
  \]
  These four sets are disjoint. If $k=|K|$, then $|L|=|M|=k>0$ and
  $\deg(u)=\deg(v)=2k$. They admit labels
  $L=\{x':x\in K\}$ and $M=\{x'':x\in K\}$ with the following complete edge rules:
  $K,L,M$ are independent; $x\sim y'$ and $x\sim y''$ exactly when $x=y$;
  $x'\sim y''$ exactly when $x\ne y$.
  For $w\in N$, $w\sim x'$ and $w\sim x''$ hold exactly when $w\not\sim x$.}
 \begin{proof}
  \step{210pairing}{For $p\in\{u,v\}$ and the other vertex $q$, the triangles through $p$
   partition $\Gamma_1(p)$ into pairs; $q$ sees exactly one vertex in each pair.}
  \begin{proof}
   \step{210involution}{On $\Gamma_1(p)$, the map $y\mapsto f(p,y)$ is a
    fixed-point-free involution, and these pairs are all the edges within $\Gamma_1(p)$.}
   \begin{proof}
    \pf\ The proof of \stepref{210triangles} gives existence and uniqueness of the third
    vertex. Thus $f(p,f(p,y))=y$. If two neighbours of $p$ are adjacent, uniqueness of
    the common neighbour on an edge through $p$ makes them one of these pairs.~\qed
   \end{proof}
   \qedstep
   \begin{proof}
    \pf\ Since $p\not\sim q$, the triangle property forces $q$ to see exactly one of
    the two other vertices of each triangle through $p$.~\qed
   \end{proof}
  \end{proof}
  \step{210labels}{For $x\in K$, put $x'=f(u,x)$ and $x''=f(v,x)$.
   These define bijections $K\to L$ and $K\to M$, and give the stated rules inside
   $K\cup L$ and $K\cup M$.}
  \begin{proof}
   \step{210onto}{Each pair through $u$ consists of one point of $K$ and one of $L$;
    each pair through $v$ consists of one point of $K$ and one of $M$.}
   \begin{proof}\pf\ Apply \stepref{210pairing} with $(p,q)=(u,v)$ and $(v,u)$, respectively.~\qed\end{proof}
   \qedstep
   \begin{proof}
    \pf\ Step \stepref{210onto} gives the bijections. Since \stepref{210pairing} lists
    every edge in either neighbourhood, each part is independent and the only edges
    between its two parts are the labelled pairs.~\qed
   \end{proof}
  \end{proof}
  \step{210lm}{$x'\sim y''$ if and only if $x\ne y$.}
  \begin{proof}
   \step{210same}{For $x=y$, $x'\not\sim x''$.}
   \begin{proof}
    \pf\ The vertex $x''$ sees $x$ and does not see $u$. Apply the triangle property to
    $\{u,x,x'\}$ to exclude $x''\sim x'$.~\qed
   \end{proof}
   \step{210different}{For $x\ne y$, $x'\sim y''$.}
   \begin{proof}
    \pf\ The vertex $x'$ sees neither $v$ nor $y$, by its definition and
    \stepref{210labels}. Applied to $\{v,y,y''\}$, the triangle property forces
    $x'\sim y''$.~\qed
   \end{proof}
   \qedstep
   \begin{proof}\pf\ The cases \stepref{210same} and \stepref{210different} are exhaustive.~\qed\end{proof}
  \end{proof}
  \step{210signature}{For $w\in N$ and $x\in K$, the adjacencies to $(x,x',x'')$
   are either $(1,0,0)$ or $(0,1,1)$.}
  \begin{proof}
   \pf\ The vertex $w$ sees neither $u$ nor $v$. It must see one vertex in each
   of $\{u,x,x'\}$ and $\{v,x,x''\}$ by the triangle property.~\qed
  \end{proof}
  \qedstep
  \begin{proof}
   \pf\ The sets partition $V(X)\diff\{u,v\}$ by their two adjacency indicators.
   Step \stepref{210labels} gives the equal cardinalities and degrees; positivity follows
   from \stepref{210triangles}. The remaining assertions are
   \stepref{210labels}, \stepref{210lm}, and \stepref{210signature}.~\qed
  \end{proof}
 \end{proof}

 \step{210regular}{\repair\ $X$ is regular. We fill the paper's missing propagation
  from nonadjacent pairs to arbitrary pairs (Gap 3 of \path{cgss/errata.md}).}
 \begin{proof}
  \step{210nonadjdeg}{Distinct nonadjacent vertices have equal degrees.}
  \begin{proof}\pf\ This is the degree equality in \stepref{210local}, for an arbitrary pair.~\qed\end{proof}
  \step{210adjdeg}{\assume{$p\sim q$.}\prove{$\deg(p)=\deg(q)$.}}
  \begin{proof}
   \pflet{$r=f(p,q)$. Choose $s\ne p$ nonadjacent to $p$, using \stepref{210triangles}.}
   \step{210s}{The vertex $s$ is outside $\{p,q,r\}$ and sees exactly one of $q,r$.}
   \begin{proof}
    \pf\ Both $q,r$ see $p$, so $s\ne q,r$; apply the triangle property and $s\not\sim p$.~\qed
   \end{proof}
   \step{210seesr}{\case{$s\sim r$.}}
   \begin{proof}
    \pf\ Step \stepref{210s} gives $s\not\sim q$. Thus
    $\deg(p)=\deg(s)=\deg(q)$ by \stepref{210nonadjdeg}.~\qed
   \end{proof}
   \step{210seesq}{\case{$s\sim q$.}}
   \begin{proof}
    \step{210prdeg}{$\deg(p)=\deg(r)$.}
    \begin{proof}
     \pf\ Step \stepref{210s} gives $s\not\sim r$;
     apply \stepref{210nonadjdeg} to $p,s$ and $s,r$.~\qed
    \end{proof}
    \pflet{Choose $t\ne q$ nonadjacent to $q$, using \stepref{210triangles}.}
    \step{210t}{The vertex $t$ is outside $\{p,q,r\}$ and sees exactly one of $p,r$.}
    \begin{proof}\pf\ Since $p,r$ see $q$, the vertex $t$ differs from both. The triangle property gives one neighbour in $\{p,q,r\}$; it cannot be $q$.~\qed\end{proof}
    \step{210tdegrees}{$\deg(q)$ equals $\deg(r)$ or $\deg(p)$.}
    \begin{proof}
     \pf\ If $t\sim p$, then $t\not\sim r$, so $\deg(q)=\deg(t)=\deg(r)$.
     If $t\sim r$, then $t\not\sim p$, so $\deg(q)=\deg(t)=\deg(p)$.
     Both equalities use \stepref{210nonadjdeg}; \stepref{210t} exhausts the cases.~\qed
    \end{proof}
    \qedstep
    \begin{proof}\pf\ Use \stepref{210prdeg} in either alternative of \stepref{210tdegrees}.~\qed\end{proof}
   \end{proof}
   \qedstep
   \begin{proof}\pf\ Step \stepref{210s} makes \stepref{210seesr} and \stepref{210seesq} exhaustive.~\qed\end{proof}
  \end{proof}
  \qedstep
  \begin{proof}
   \pf\ Identical vertices have equal degrees; distinct vertices are covered by
   \stepref{210nonadjdeg} and \stepref{210adjdeg}.~\qed
  \end{proof}
 \end{proof}

 \pflet{Fix arbitrary distinct nonadjacent $u,v$ and retain the sets and labels of
 \stepref{210local}. Put $k=|K|$. By \stepref{210regular}, every vertex has degree $2k$.}
 \step{210kge}{$k\geq2$.}
 \begin{proof}
  \step{210kone}{If $k=1$, then $X$ is a single triangle.}
  \begin{proof}
   \pf\ By \stepref{210regular}, every vertex would have degree 2. A triangle exists by
   \stepref{210triangles}; its vertices already have two neighbours inside it, whereas
   every outside vertex would have a neighbour in it. Thus there is no outside vertex.~\qed
  \end{proof}
  \qedstep
  \begin{proof}
   \pf\ Step \stepref{210local} gives $k>0$; \stepref{210kone} contradicts the chosen nonadjacent pair when $k=1$.~\qed
  \end{proof}
 \end{proof}

 \step{210counts}{$|N|=3k-5$, $|V(X)|=6k-3$, and $X$ has strongly regular parameters
  $(6k-3,2k,1,k)$. The graph on $N$ has $k(k-2)$ edges, and the third vertex
  of every such edge lies in $K$.}
 \begin{proof}
  \step{210cross}{There are $4k^2-6k$ edges between $K\cup L\cup M$ and $N$.}
  \begin{proof}
   \step{210degreesintoN}{A vertex in $K$ has $2k-4$ neighbours in $N$;
    a vertex in $L$ or $M$ has $k-1$ neighbours in $N$.}
   \begin{proof}
    \pf\ By \stepref{210local}, $x\in K$ has the four neighbours $u,v,x',x''$
    outside $N$. A vertex $x'\in L$ has $u,x$ and the $k-1$ vertices $y''$ with
    $y\ne x$ outside $N$; a vertex $x''\in M$ has $v,x$ and the $k-1$ vertices
    $y'$ with $y\ne x$. Subtract these counts from degree $2k$.~\qed
   \end{proof}
   \qedstep
   \begin{proof}
    \pf\ Sum \stepref{210degreesintoN} over the three parts, each of size $k$:
    $k(2k-4)+k(k-1)+k(k-1)=4k^2-6k$.~\qed
   \end{proof}
  \end{proof}
  \step{210Nedges}{The subgraph on $N$ has $k(k-2)$ edges, each with its third vertex in $K$.}
  \begin{proof}
   \step{210Nthird}{If $w,z\in N$ are adjacent, then $f(w,z)\in K$.}
   \begin{proof}
    \step{210Noutside}{Both $u$ and $v$ lie outside $\{w,z,f(w,z)\}$.}
    \begin{proof}
     \pf\ They differ from $w,z$ by the definition of $N$.
     The vertex $f(w,z)$ is adjacent to $w$, whereas neither $u$ nor $v$ is;
     thus they also differ from $f(w,z)$.~\qed
    \end{proof}
    \qedstep
    \begin{proof}
     \pf\ By \stepref{210Noutside}, the triangle property applies to $u$ and $v$.
     Each sees neither $w$ nor $z$, so each sees $f(w,z)$. This is membership in $K$.~\qed
    \end{proof}
   \end{proof}
   \step{210Kpairs}{Each $x\in K$ is the third vertex of exactly $k-2$ edges in $N$.}
   \begin{proof}
    \pf\ In the triangle pairing of the $2k$ neighbours of $x$, the pairs
    $\{u,x'\}$ and $\{v,x''\}$ use its four neighbours outside $N$.
    The remaining $2k-4$ neighbours pair within $N$, giving $k-2$ edges.~\qed
   \end{proof}
   \qedstep
   \begin{proof}
    \pf\ Step \stepref{210Nthird} covers every edge in $N$; uniqueness of its third
    vertex prevents double counting. Sum \stepref{210Kpairs} over $|K|=k$.~\qed
   \end{proof}
  \end{proof}
  \step{210cardN}{$|N|=3k-5$ and $|V(X)|=6k-3$.}
  \begin{proof}
   \pf\ Counting the degree sum over $N$, using \stepref{210cross} and \stepref{210Nedges},
   gives
   \begin{align*}
    2k|N|&=2k(k-2)+(4k^2-6k)&&\text{[internal edges count twice]}\\
          &=2k(3k-5)&&\text{[ring arithmetic],}\\
    |N|&=3k-5&&\text{[$2k>0$],}\\
    |V(X)|&=2+3k+(3k-5)=6k-3&&\text{[the partition].}
   \end{align*}
   ~\qed
  \end{proof}
  \qedstep
  \begin{proof}
   \pf\ Step \stepref{210cardN} gives the order and \stepref{210regular} the degree.
   Step \stepref{210triangles} gives one common neighbour on each edge.
   For any nonadjacent pair, \stepref{210local} gives half its degree as the number
   of common neighbours, namely $k$. These are the four strongly regular parameters.~\qed
  \end{proof}
 \end{proof}

 \step{210dK}{\repair\ For $w\in N$, put
  $A(w)=K\diff\Gamma_1(w)$ and $d_K(w)=|K\cap\Gamma_1(w)|$.
  Then $|A(w)|\in\{0,2\}$, $d_K(w)\in\{k-2,k\}$, and
  $|N\cap\Gamma_1(w)|=d_K(w)$. We replace the unsupported pigeonhole argument
  by this invariant (Gap 1 of \path{cgss/errata.md}).}
 \begin{proof}
  \step{210KNpairing}{In the triangle pairing at $w$, its neighbours in $K$ pair
   with its neighbours in $N$. Its neighbours in $L$ pair with its neighbours in $M$.}
  \begin{proof}
   \step{210KtoN}{If $w\sim x\in K$, then $f(w,x)\in N$.}
   \begin{proof}
    \pf\ Each of $u,v$ is adjacent to $x$ and not to $w$. The triangle property
    excludes adjacency to $f(w,x)$. Also $f(w,x)\ne u,v$, since both are nonadjacent
    to $w$. The definition of $N$ applies.~\qed
   \end{proof}
   \step{210NtoK}{If $w\sim z\in N$, then $f(w,z)\in K$.}
   \begin{proof}
    \pf\ Apply the third-vertex assertion of \stepref{210counts} to the edge $wz$.~\qed
   \end{proof}
   \step{210LMclosed}{The partner of a neighbour of $w$ in $L\cup M$ lies in $L\cup M$.}
   \begin{proof}
    \pflet{Let $y\in(L\cup M)\cap\Gamma_1(w)$ and put $z=f(w,y)$.}
    \step{210partnerback}{$f(w,z)=y$.}
    \begin{proof}\pf\ The vertex $y$ is a common neighbour of $w,z$;
     uniqueness in \stepref{210triangles} gives the equality.~\qed\end{proof}
    \step{210partnernotUV}{$z\notin\{u,v\}$.}
    \begin{proof}\pf\ The vertex $z$ is adjacent to $w$, whereas $u,v$ are not.~\qed\end{proof}
    \step{210partnernotK}{$z\notin K$.}
    \begin{proof}\pf\ If $z\in K$, \stepref{210KtoN} gives
     $f(w,z)\in N$, contrary to \stepref{210partnerback} and $y\in L\cup M$.~\qed\end{proof}
    \step{210partnernotN}{$z\notin N$.}
    \begin{proof}\pf\ If $z\in N$, \stepref{210NtoK} gives
     $f(w,z)\in K$, contrary to \stepref{210partnerback} and $y\in L\cup M$.~\qed\end{proof}
    \qedstep
    \begin{proof}\pf\ The partition in \stepref{210local} and
     \stepref{210partnernotUV}, \stepref{210partnernotK}, \stepref{210partnernotN}
     leave only $L\cup M$.~\qed\end{proof}
   \end{proof}
   \qedstep
   \begin{proof}
    \pf\ The involution $y\mapsto f(w,y)$ pairs the neighbours of $w$.
    Steps \stepref{210KtoN} and \stepref{210NtoK} pair all its $K$ and $N$ neighbours.
    Step \stepref{210LMclosed} keeps every remaining pair in $L\cup M$.
    Neither part has internal edges by \stepref{210local}, so each such pair has
    one vertex in each part.~\qed
   \end{proof}
  \end{proof}
  \step{210Asize}{If $A(w)\ne\varnothing$, then $|A(w)|=2$.}
  \begin{proof}
   \pflet{Choose $x\in A(w)$.}
   \step{210common}{The common neighbours of $w,x'$ are exactly
    $\{y'':y\in A(w),\ y\ne x\}$.}
   \begin{proof}
    \pf\ The signature in \stepref{210local} gives $w\sim x'$.
    By \stepref{210KNpairing} their unique common neighbour lies in $M$.
    For $y\in K$, we have $w\sim y''$ exactly when $y\in A(w)$, and
    $x'\sim y''$ exactly when $y\ne x$, by \stepref{210local}.
    Thus the displayed set lists all, and only, common neighbours.~\qed
   \end{proof}
   \qedstep
   \begin{proof}
    \pf\ An edge has exactly one common neighbour by \stepref{210triangles}.
    The bijection $y\mapsto y''$ in \stepref{210local} and \stepref{210common}
    give $|A(w)|-1=1$.~\qed
   \end{proof}
  \end{proof}
  \qedstep
  \begin{proof}
   \pf\ Either $A(w)$ is empty or \stepref{210Asize} applies.
   Since $d_K(w)=k-|A(w)|$, its two possible values follow.
   The $K,N$ pairing in \stepref{210KNpairing} gives the asserted equality of degrees.~\qed
  \end{proof}
 \end{proof}

 \step{210degreeSum}{$\sum_{w\in N}d_K(w)=k(2k-4)$.}
 \begin{proof}
  \pf\ Each $x\in K$ has $2k-4$ neighbours in $N$: its four neighbours outside $N$
  are $u,v,x',x''$ by \stepref{210local}, and its degree is $2k$.
  Count edges between $K,N$ from the two ends.~\qed
 \end{proof}
 \pflet{$h\triangleq|\{w\in N:d_K(w)=k\}|$.}
 \step{210highcount}{$2h=(k-2)(5-k)$.}
 \begin{proof}
  \pf\ By \stepref{210dK} and \stepref{210counts}, the sum in \stepref{210degreeSum} is
  $(|N|-h)(k-2)+hk=(3k-5)(k-2)+2h$.
  Equating the two expressions and subtracting gives the formula.~\qed
 \end{proof}
 \step{210combbound}{\new\ $2\leq k\leq5$. We obtain the upper bound without a
  Gram matrix, so the result applies to an arbitrary graph with the triangle property.}
 \begin{proof}
  \step{210negative}{If $k>5$, then $(k-2)(5-k)<0$.}
  \begin{proof}\pf\ We have $k-2>0$ and $5-k<0$; multiply these inequalities.~\qed\end{proof}
  \qedstep
  \begin{proof}
   \pf\ The cardinality $h$ is nonnegative. Thus \stepref{210highcount} contradicts
   \stepref{210negative} if $k>5$. The lower bound is \stepref{210kge}.~\qed
  \end{proof}
 \end{proof}
 \step{210exclude}{\repair\ We have $k\in\{2,3,5\}$ (Gap 1).}
 \begin{proof}
  \step{210highforces}{If $h>0$, then $k=3$.}
  \begin{proof}
   \pflet{Choose $w_0\in N$ with $d_K(w_0)=k$.}
   \step{210choosez}{There is $z\in N$ adjacent to $w_0$.}
   \begin{proof}\pf\ Steps \stepref{210dK} and \stepref{210kge} give at least two neighbours of $w_0$ in $N$.~\qed\end{proof}
   \pflet{Choose such a $z$ and put $x=f(w_0,z)$.}
   \step{210zone}{$K\cap\Gamma_1(z)=\{x\}$.}
   \begin{proof}
    \pf\ Step \stepref{210counts} puts $x$ in $K$.
    If $y\in K\diff\{x\}$, then $y\sim w_0$ and $y\not\sim x$, since $K$ is independent.
    The triangle property for $\{w_0,z,x\}$ forces $y\not\sim z$.~\qed
   \end{proof}
   \qedstep
   \begin{proof}
    \pf\ Step \stepref{210zone} gives $d_K(z)=1$.
    By \stepref{210dK}, $1=k-2$ or $1=k$; \stepref{210kge} excludes the second,
    so $k=3$.~\qed
   \end{proof}
  \end{proof}
  \qedstep
  \begin{proof}
   \pf\ Step \stepref{210combbound} leaves $k=2,3,4,5$.
   At $k=4$, \stepref{210highcount} gives $h=1$, whereas
   \stepref{210highforces} gives $k=3$. Thus $k=4$ is impossible.~\qed
  \end{proof}
 \end{proof}

 \step{210two}{\repair\ (Gap 2.) If $k=2$, the labelled graph is uniquely determined:
  $N$ consists of one vertex $w$, with $A(w)=K$, and has no edges.}
 \begin{proof}
  \step{210twoN}{$|N|=1$ and $d_K(w)=0$.}
  \begin{proof}
   \pf\ Step \stepref{210counts} gives $|N|=1$.
   A simple graph on one vertex has degree 0, so \stepref{210dK} gives $d_K(w)=0$.~\qed
  \end{proof}
  \qedstep
  \begin{proof}
   \pf\ Step \stepref{210twoN} gives $A(w)=K$.
   The signature in \stepref{210local} makes $w$ adjacent to every vertex of $L\cup M$
   and to none of $K\cup\{u,v\}$. The remaining edges are all listed in
   \stepref{210local}. Thus every edge is determined by labels on the two elements of $K$.~\qed
  \end{proof}
 \end{proof}

 \step{210three}{\repair\ (Gap 2.) If $k=3$, the subgraph on $N$ is a claw. Its centre sees all
  of $K$ and its three leaves correspond bijectively to $K$ by their unique neighbour there.
  These labels determine the whole graph.}
 \begin{proof}
  \step{210centre}{There is exactly one $w_0\in N$ with $d_K(w_0)=3$;
   the other three vertices have $d_K=1$.}
  \begin{proof}
   \pf\ Step \stepref{210highcount} gives $2h=(3-2)(5-3)=2$.
   Thus $h=1$; \stepref{210counts} gives $|N|=4$, and \stepref{210dK}
   gives the values 1 or 3.~\qed
  \end{proof}
  \step{210claw}{The vertex $w_0$ is adjacent to all three other vertices of $N$,
   and they are pairwise nonadjacent.}
  \begin{proof}
   \pf\ By \stepref{210dK}, the degrees inside $N$ are 3 for $w_0$ and 1 for
   the others. Its three edges use the entire degree of each of the other vertices.~\qed
  \end{proof}
  \step{210leaflabels}{For each $x\in K$ there is exactly one leaf $w_x$ with
   $K\cap\Gamma_1(w_x)=\{x\}$.}
  \begin{proof}
   \pf\ By \stepref{210local} and \stepref{210regular}, the vertex $x$ has
   $2k-4=2$ neighbours in $N$, of which one is $w_0$.
   Its other neighbour is a leaf. Conversely, each leaf has exactly one neighbour in $K$
   by \stepref{210centre}. Thus this incidence is a bijection.~\qed
  \end{proof}
  \qedstep
  \begin{proof}
   \pf\ Steps \stepref{210centre}, \stepref{210claw}, and \stepref{210leaflabels}
   determine $N$ and its incidences with $K$.
   The signature in \stepref{210local} determines its incidences with $L,M$;
   the same step already lists all other edges.~\qed
  \end{proof}
 \end{proof}

 \step{210five}{\repair\ (Gap 2.) If $k=5$, the map
  $w\mapsto A(w)=K\diff\Gamma_1(w)$ is a bijection from $N$ to the two-element subsets
  of $K$. Two vertices of $N$ are adjacent exactly when their labels are disjoint.
  These labels determine the whole graph.}
 \begin{proof}
  \step{210pairsizes}{$d_K(w)=3$ and $|A(w)|=2$ for every $w\in N$.}
  \begin{proof}
   \pf\ Step \stepref{210counts} gives $|N|=10$ and \stepref{210degreeSum} gives
   $\sum_{w\in N}d_K(w)=5(10-4)=30$. Each summand is 3 or 5 by \stepref{210dK};
   the minimum total is already 30, so every summand is 3.~\qed
  \end{proof}
  \step{210pairlabel}{For distinct $x,y\in K$, there is exactly one $w\in N$
   with $A(w)=\{x,y\}$.}
  \begin{proof}
   \step{210edgeLM}{$x'$ and $y''$ are adjacent and their third vertex lies in $N$.}
   \begin{proof}
    \step{210noCore}{They have no common neighbour in $\{u,v\}\cup K\cup L\cup M$.}
    \begin{proof}
     \pf\ The vertex $u$ does not see $y''$, and $v$ does not see $x'$.
     Their respective neighbours in $K$ are the distinct vertices $x,y$.
     The set $L$ is independent and contains $x'$; the set $M$ is independent and
     contains $y''$. These observations exclude each of the five parts.~\qed
    \end{proof}
    \qedstep
    \begin{proof}
     \pf\ Step \stepref{210local} gives $x'\sim y''$.
     Their unique common neighbour exists by \stepref{210triangles};
     \stepref{210noCore} and the partition put it in $N$.~\qed
    \end{proof}
   \end{proof}
   \step{210labelvalue}{For $w=f(x',y'')$, $A(w)=\{x,y\}$.}
   \begin{proof}
    \pf\ Since $w$ sees $x',y''$, the signature in \stepref{210local} gives
    $x,y\in A(w)$. Step \stepref{210pairsizes} gives $|A(w)|=2$, so these are all its elements.~\qed
   \end{proof}
   \step{210labelunique}{Any $z\in N$ with $A(z)=\{x,y\}$ equals $w$.}
   \begin{proof}
    \pf\ The signature makes $z$ a common neighbour of $x',y''$;
    uniqueness of the third vertex gives $z=f(x',y'')=w$.~\qed
   \end{proof}
   \qedstep
   \begin{proof}\pf\ Steps \stepref{210labelvalue} and \stepref{210labelunique} give existence and uniqueness.~\qed\end{proof}
  \end{proof}
  \step{210disjoint}{For $w,z\in N$, $w\sim z$ if and only if $A(w)\cap A(z)=\varnothing$.}
  \begin{proof}
   \step{210adjdisjoint}{If $w\sim z$, then $A(w)\cap A(z)=\varnothing$.}
   \begin{proof}
    \pf\ Their third vertex $x$ is in $K$ by \stepref{210counts}.
    Every $y\in K\diff\{x\}$ is
    nonadjacent to $x$, so the triangle property makes it adjacent to exactly one of
    $w,z$; $x$ is adjacent to both. Thus no element of $K$ is missed by both.~\qed
   \end{proof}
   \step{210threechoices}{For fixed $w$, there are exactly three $z\in N$ whose labels
    are disjoint from $A(w)$, and exactly three neighbours of $w$ in $N$.}
   \begin{proof}
    \pf\ The complement of $A(w)$ in $K$ has three elements, and thus three two-element
    subsets. Step \stepref{210pairlabel} gives exactly one vertex for each.
    Steps \stepref{210dK} and \stepref{210pairsizes} give degree 3 within $N$.~\qed
   \end{proof}
   \qedstep
   \begin{proof}
    \pf\ Step \stepref{210adjdisjoint} gives inclusion of the neighbour set in the
    disjoint-label set; \stepref{210threechoices} gives equal finite cardinalities,
    so the two sets coincide.~\qed
   \end{proof}
  \end{proof}
  \qedstep
  \begin{proof}
   \pf\ Steps \stepref{210pairsizes} and \stepref{210pairlabel} give the bijection,
   and \stepref{210disjoint} gives the graph on $N$, the two-subset model of the Petersen
   graph. The labels give all incidences with $K$, their complements give those with
   $L,M$ by \stepref{210local}, and there are no incidences with $u,v$.~\qed
  \end{proof}
 \end{proof}

 \qedstep
 \begin{proof}
  \pf\ Steps \stepref{210local}, \stepref{210regular}, and \stepref{210counts}
  give the partition, edge rules, regularity and parameters.
  Steps \stepref{210dK}, \stepref{210degreeSum}, and \stepref{210highcount}
  give the degree identities. Step \stepref{210exclude} leaves exactly the three cases
  proved in \stepref{210two}, \stepref{210three}, and \stepref{210five}.
  The pair $u,v$ was arbitrary.~\qed
 \end{proof}
\end{proof}

\begin{namedresult}{Lemma}{(isomorphism from the order)}\label{lem:triangle-iso}
 \lean{SimpleGraph.HasTriangleProperty.nonempty_iso_of_card_eq}
 \uses{def:triangle-property,lem:triangle-normal-form}
 \assume{\begin{pfenum}
  \item $X,Y$ are finite simple graphs with the triangle property, and each has an edge.
  \item Neither graph has a vertex adjacent to every other vertex.
  \item $|V(X)|=|V(Y)|$.
 \end{pfenum}}
 \prove{There is a graph isomorphism $X\to Y$.}
\end{namedresult}
\begin{proof}
 \pfsketch\ \repair\ We will apply the normal-form lemma separately to the two graphs
 and match their labels. This completes the uniqueness argument of Gap 2.
 \step{210pairs}{Choose distinct nonadjacent $u,v\in V(X)$ and
  $\widetilde u,\widetilde v\in V(Y)$.}
 \begin{proof}
  \pf\ Each graph has a vertex since it has an edge. By \asm[2]{0}, that vertex
  has a distinct non-neighbour, in each graph.~\qed
 \end{proof}
 \step{210formX}{The normal-form lemma applies to $X,u,v$; write $k_X=|K|$.
  In particular $|V(X)|=6k_X-3$.}
 \begin{proof}\pf\ Apply Lemma~\ref{lem:triangle-normal-form} to the hypotheses on $X$
  and the first pair in \stepref{210pairs}.~\qed\end{proof}
 \step{210formY}{The normal-form lemma applies to $Y,\widetilde u,\widetilde v$;
  write $k_Y=|\widetilde K|$. In particular $|V(Y)|=6k_Y-3$.}
 \begin{proof}\pf\ Apply Lemma~\ref{lem:triangle-normal-form} to the hypotheses on $Y$
  and the second pair in \stepref{210pairs}.~\qed\end{proof}
 \step{210samek}{$k_X=k_Y$.}
 \begin{proof}\pf\ By \asm[3]{0}, \stepref{210formX} and \stepref{210formY},
  $6k_X-3=6k_Y-3$; cancellation gives the equality.~\qed\end{proof}
 \pflet{Write $k$ for this common value, use tildes for the parts of $Y$,
  and choose a bijection $\beta:K\to\widetilde K$, using \stepref{210samek}.}
 \step{210iso}{\repair\ The bijection $\beta$ extends to a graph isomorphism (Gap 2).}
 \begin{proof}
  \step{210extend}{Extend $\beta$ by $u\mapsto\widetilde u$, $v\mapsto\widetilde v$,
   $x'\mapsto\beta(x)'$, $x''\mapsto\beta(x)''$, and the following rule on $N$:
   send the single vertex to the single vertex when $k=2$; send centre to centre and
   $w_x$ to $\widetilde w_{\beta(x)}$ when $k=3$; send label $A$ to label $\beta(A)$
   when $k=5$. The resulting map $\psi$ is a bijection.}
  \begin{proof}
   \step{210extendcore}{The prescribed maps on $\{u\},\{v\},K,L,M$ are bijections
    onto the corresponding parts of $Y$.}
   \begin{proof}\pf\ The label maps from $K$ to $L,M$ are bijections in both
    \stepref{210formX} and \stepref{210formY}; compose these with $\beta$.~\qed\end{proof}
   \step{210extendtwo}{When $k=2$, the map on $N$ is a bijection.}
   \begin{proof}\pf\ Both sets are singletons by part 6 of the normal-form lemma.~\qed\end{proof}
   \step{210extendthree}{When $k=3$, the map on $N$ is a bijection.}
   \begin{proof}\pf\ Part 7 of the normal-form lemma gives one centre and a
    bijective leaf labelling by $K$ in each graph; $\beta$ matches the labels.~\qed\end{proof}
   \step{210extendfive}{When $k=5$, the map on $N$ is a bijection.}
   \begin{proof}\pf\ Part 8 of the normal-form lemma identifies both sets with
    two-element subsets, and $A\mapsto\beta(A)$ is a bijection on these subsets.~\qed\end{proof}
   \qedstep
   \begin{proof}\pf\ The parts partition both vertex sets by the normal-form lemma.
    Step \stepref{210extendcore} and the exhaustive cases
    \stepref{210extendtwo}, \stepref{210extendthree}, \stepref{210extendfive}
    give a bijection on their disjoint unions.~\qed\end{proof}
  \end{proof}
  \step{210preserve}{The map $\psi$ preserves adjacency and nonadjacency.}
  \begin{proof}
   \step{210coreedges}{It preserves all pairs with neither endpoint in $N$.}
   \begin{proof}
    \step{210coreuv}{The pair $u,v$ is a nonedge in each graph.}
    \begin{proof}\pf\ This is the choice in \stepref{210pairs}.~\qed\end{proof}
    \step{210coreuK}{Every pair $u,x$ with $x\in K$ is an edge in each graph.}
    \begin{proof}\pf\ This is the definition of $K$.~\qed\end{proof}
    \step{210coreuL}{Every pair $u,x'$ with $x'\in L$ is an edge in each graph.}
    \begin{proof}\pf\ This is the definition of $L$.~\qed\end{proof}
    \step{210coreuM}{Every pair $u,x''$ with $x''\in M$ is a nonedge in each graph.}
    \begin{proof}\pf\ This is the definition of $M$.~\qed\end{proof}
    \step{210corevK}{Every pair $v,x$ with $x\in K$ is an edge in each graph.}
    \begin{proof}\pf\ This is the definition of $K$.~\qed\end{proof}
    \step{210corevL}{Every pair $v,x'$ with $x'\in L$ is a nonedge in each graph.}
    \begin{proof}\pf\ This is the definition of $L$.~\qed\end{proof}
    \step{210corevM}{Every pair $v,x''$ with $x''\in M$ is an edge in each graph.}
    \begin{proof}\pf\ This is the definition of $M$.~\qed\end{proof}
    \step{210coresamePart}{Pairs within any one part $D\in\{K,L,M\}$ are nonedges in each graph.}
    \begin{proof}\pf\ Each such $D$ is independent by part 3 of the normal-form lemma.~\qed\end{proof}
    \step{210coreKL}{The relation $x\sim y'$ is preserved in both directions.}
    \begin{proof}\pf\ In each graph it is $x=y$, by part 3 of the normal-form lemma; $\beta$ preserves and reflects equality.~\qed\end{proof}
    \step{210coreKM}{The relation $x\sim y''$ is preserved in both directions.}
    \begin{proof}\pf\ In each graph it is $x=y$, by part 3 of the normal-form lemma; $\beta$ preserves and reflects equality.~\qed\end{proof}
    \step{210coreLM}{The relation $x'\sim y''$ is preserved in both directions.}
    \begin{proof}\pf\ In each graph it is $x\ne y$, by part 3 of the normal-form lemma; $\beta$ preserves and reflects inequality.~\qed\end{proof}
    \qedstep
    \begin{proof}\pf\ The listed pair types exhaust the five parts
     $\{u\},\{v\},K,L,M$. Loops are nonedges because the graphs are simple.~\qed\end{proof}
   \end{proof}
   \step{210mixededges}{It preserves all pairs with exactly one endpoint in $N$.}
   \begin{proof}
    \step{210mixeduv}{Vertices of $N$ are adjacent to neither $u$ nor $v$ in either graph.}
    \begin{proof}\pf\ This is the definition of $N$.~\qed\end{proof}
    \step{210mixedtwo}{For $k=2$, $\beta(A(w))=\widetilde A(\psi(w))$.}
    \begin{proof}\pf\ Both missed sets are the whole of $K$ or $\widetilde K$,
     respectively, by part 6 of the normal-form lemma.~\qed\end{proof}
    \step{210mixedthree}{For $k=3$, $\beta(A(w))=\widetilde A(\psi(w))$.}
    \begin{proof}
     \step{210mixedcentre}{The equality holds at the centre.}
     \begin{proof}\pf\ Both missed sets are empty by part 7 of the normal-form lemma.~\qed\end{proof}
     \step{210mixedleaf}{The equality holds at each leaf $w_x$.}
     \begin{proof}\pf\ Its missed set is $K\diff\{x\}$ and its image has missed set
      $\widetilde K\diff\{\beta(x)\}$, by part 7 of the normal-form lemma.~\qed\end{proof}
     \qedstep
     \begin{proof}\pf\ A claw consists of its centre and leaves; apply
      \stepref{210mixedcentre} and \stepref{210mixedleaf}.~\qed\end{proof}
    \end{proof}
    \step{210mixedfive}{For $k=5$, $\beta(A(w))=\widetilde A(\psi(w))$.}
    \begin{proof}\pf\ This is the definition of $\psi$ on the two-subset labels
     in \stepref{210extend}.~\qed\end{proof}
    \step{210mixedK}{Pairs between $N$ and $K$ are preserved in both directions.}
    \begin{proof}\pf\ Adjacency is $x\notin A(w)$ by definition.
     Steps \stepref{210mixedtwo}, \stepref{210mixedthree}, and \stepref{210mixedfive}
     preserve this condition in every possible case.~\qed\end{proof}
    \step{210mixedL}{Pairs between $N$ and $L$ are preserved in both directions.}
    \begin{proof}\pf\ The relation $w\sim x'$ is $x\in A(w)$ by part 4 of the
     normal-form lemma, preserved by the three missed-set equalities above.~\qed\end{proof}
    \step{210mixedM}{Pairs between $N$ and $M$ are preserved in both directions.}
    \begin{proof}\pf\ The relation $w\sim x''$ is $x\in A(w)$ by part 4 of the
     normal-form lemma, preserved by the same equalities.~\qed\end{proof}
    \qedstep
    \begin{proof}\pf\ Steps \stepref{210mixeduv}, \stepref{210mixedK},
     \stepref{210mixedL}, and \stepref{210mixedM} exhaust the other endpoint.~\qed\end{proof}
   \end{proof}
   \step{210internaledges}{It preserves all pairs with both endpoints in $N$.}
   \begin{proof}
    \step{210internaltwo}{When $k=2$, there are no edges in either $N$.}
    \begin{proof}\pf\ Both sets are singletons by part 6 of the normal-form lemma.~\qed\end{proof}
    \step{210internalthree}{When $k=3$, adjacency within $N$ is preserved in both directions.}
    \begin{proof}\pf\ Part 7 of the normal-form lemma says that the edges join
     the centre to a leaf, and \stepref{210extend} preserves these roles.~\qed\end{proof}
    \step{210internalfive}{When $k=5$, adjacency within $N$ is preserved in both directions.}
    \begin{proof}\pf\ Part 8 of the normal-form lemma says that adjacency is disjointness
     of labels; $\beta(A)\cap\beta(B)=\beta(A\cap B)$ since $\beta$ is injective.~\qed\end{proof}
    \qedstep
    \begin{proof}\pf\ The normal-form lemma gives $k\in\{2,3,5\}$, so the three
     preceding cases are exhaustive.~\qed\end{proof}
   \end{proof}
   \qedstep
   \begin{proof}\pf\ Steps \stepref{210coreedges}, \stepref{210mixededges}, and
    \stepref{210internaledges} exhaust all pairs.~\qed\end{proof}
  \end{proof}
  \qedstep
  \begin{proof}\pf\ The bijection \stepref{210extend} and the equivalence of
   adjacency in \stepref{210preserve} define a graph isomorphism.~\qed\end{proof}
 \end{proof}
 \qedstep
 \begin{proof}\pf\ Step \stepref{210iso} gives the required isomorphism.~\qed\end{proof}
\end{proof}

\begin{cgss}{Lemma}{2.10}\label{lem:2.10}
 \lean{LineSystem.Hyp24.card_Γa_mem}
 \lean{LineSystem.Hyp24.orthGraph_iso_of_card_Γa_eq}
 \lean{LineSystem.Hyp24.orthGraph_isSRGWith}
 \uses{hyp:2.4,lem:2.6,lem:2.9,cor:triangle-property,lem:triangle-normal-form,lem:triangle-iso,facts:inner,comp:C1}
 \assume{\begin{pfenum}
  \item Hypothesis 2.4 holds in the finite-dimensional real inner product space $V$,
   and $X$ is its orthogonality graph on $\Ga$.
  \item $X$ has an edge, is noncomplete, and no vertex is adjacent to every other vertex.
 \end{pfenum}}
 \prove{$|\Ga|\in\{9,15,27\}$, and the graph is unique up to isomorphism in each case.
  More explicitly, if another instance of Hypothesis 2.4 in a finite-dimensional real inner
  product space $W$ satisfies the same graph hypotheses and has the same value of $|\Ga|$,
  its orthogonality graph is isomorphic to $X$.
  Moreover, \repair\ $X$ is strongly regular with parameters $(6k-3,2k,1,k)$ for some
  $k\in\{2,3,5\}$ (Error 12).}
\end{cgss}
\begin{proof}
 \pfsketch\ The two graph lemmas give the conclusion. We will also retain the paper's
 Gram-matrix proof of the upper bound on $k$, including computation C1.
 By Lemma 2.9, having a universal vertex is equivalent here to being a friendship graph.
 Thus our precise hypothesis matches the paper's exclusion of that family; we retain its
 noncompleteness hypothesis as well.
 \step{210apply}{For any distinct nonadjacent $u,v$ and $k=|\Gamma_1(u)\cap\Gamma_1(v)|$,
  the normal-form conclusions hold. In particular $|\Ga|\in\{9,15,27\}$ and
  $X$ has the stated strongly regular parameters.}
 \begin{proof}
  \pf\ Corollary~\ref{cor:triangle-property} supplies the triangle property.
  Apply Lemma~\ref{lem:triangle-normal-form} and substitute $k=2,3,5$ in $6k-3$.~\qed
 \end{proof}
 \step{210unique}{Any second instance with the same graph hypotheses and order has
  an orthogonality graph isomorphic to $X$.}
 \begin{proof}
  \raggedright\pf\ Corollary~\ref{cor:triangle-property} gives the triangle property for each
  instance. Apply Lemma~\ref{lem:triangle-iso} to the two graphs.~\qed
 \end{proof}
 \step{210paperBound}{The paper's Gram-matrix argument gives a second proof of $k\leq5$.}
 \begin{proof}
  \pflet{Choose distinct nonadjacent $u,v$ and use the partition and labels of
   Lemma~\ref{lem:triangle-normal-form}. Write $\Zu,\Zv,\Zx,\Zx',\Zx''$ for their
   representing vectors and put $k=|K|\geq2$, by \stepref{210apply}.}
 \step{210gram}{The following are Gram matrices, first for the family
  \[\Zu,\Zv,(\Zx)_{x\in K},(\Zx')_{x\in K},(\Zx'')_{x\in K},\]
  then for its five block sums
  \[\Zu,\Zv,\sum_{x\in K}\Zx,\sum_{x\in K}\Zx',\sum_{x\in K}\Zx''.\]
  \[
  B_0=\begin{pmatrix}
  2&1&0^T&0^T&\Zj^T\\
  1&2&0^T&\Zj^T&0^T\\
  0&0&I+J&J-I&J-I\\
  0&\Zj&J-I&I+J&I\\
  \Zj&0&J-I&I&I+J
  \end{pmatrix},
  \]
  \[
  B=\begin{pmatrix}
  2&1&0&0&k\\
  1&2&0&k&0\\
  0&0&k^2+k&k^2-k&k^2-k\\
  0&k&k^2-k&k^2+k&k\\
  k&0&k^2-k&k&k^2+k
  \end{pmatrix}.
  \]
  Here $I,J$ have size $k$ and $\Zj$ has $k$ entries all equal to 1.
  We number the displayed rows and columns from 1.}
 \begin{proof}
  \step{210entries}{For distinct vertices $x,y$, $(\Zx,\Zy)=0$ on edges and 1 on nonedges.}
  \begin{proof}\pf\ Use the definition of $X$ and Lemma 2.6; diagonal products are 2.~\qed\end{proof}
  \step{210blocks}{The displayed blocks of $B_0$ follow from Lemma~\ref{lem:triangle-normal-form}.}
  \begin{proof}
   \pf\ Independent parts give $I+J$; matched edges give $J-I$;
   all edges except matched pairs give $I$. The products with $\Zu,\Zv$ follow from
   the definitions of $K,L,M$, using \stepref{210entries}.~\qed
  \end{proof}
  \step{210sums}{The matrix $B$ consists of the block sums of $B_0$.}
  \begin{proof}
   \pf\ Bilinearity gives each sum-vector product as the corresponding double sum.
   The sums of the entries of $I+J,J-I,I$ are $k^2+k,k^2-k,k$, respectively;
   those of $\Zj$ and 0 are $k$ and 0.~\qed
  \end{proof}
  \qedstep
  \begin{proof}
   \pf\ Steps \stepref{210blocks} and \stepref{210sums} prove the formulas.
   \repair\ In particular $B(3,1)=0$, correcting the printed 2 (Error 3 of
   \path{cgss/errata.md}): for every $x\in K$, the vector $\Zx$ is orthogonal to $\Zu$.~\qed
  \end{proof}
 \end{proof}

 \step{210bound}{$2\leq k\leq5$.}
 \begin{proof}
  \step{210psd}{The matrix $D$ obtained by deleting row and column 1 of $B$ is positive
   semidefinite, so $\det D\geq0$.}
  \begin{proof}
   \step{210Dgram}{$D$ is the Gram matrix of
    $\Zv,\sum_{x\in K}\Zx,\sum_{x\in K}\Zx',\sum_{x\in K}\Zx''$.}
   \begin{proof}\pf\ Delete the first vector from the block-sum family in
    \stepref{210gram}.~\qed\end{proof}
   \step{210Dpsd}{$D$ is positive semidefinite.}
   \begin{proof}
    \pf\ Apply \texttt{Matrix.posSemidef\_gram} in
    \path{Mathlib/Analysis/InnerProductSpace/GramMatrix.lean}
    directly to the four-vector family in \stepref{210Dgram}.~\qed
   \end{proof}
   \qedstep
   \begin{proof}
    \raggedright\pf\ Apply \path{Matrix.PosSemidef.det_nonneg} in
    \path{Mathlib/Analysis/Matrix/PosDef.lean} to \stepref{210Dpsd}.
    We index the real matrix by \texttt{Fin 4}, a finite type with decidable equality.~\qed
   \end{proof}
  \end{proof}
  \step{210det}{Computation C1: $\det D=2k^5(5-k)$.}\label{comp:C1}
  \lean{Computations.C1.det_blockSumMinor}\uses{}
  \begin{proof}
   \raggedright\pf\ Expanding along the first row $(2,0,k,0)$ gives
   \begin{align*}
    \det D&=2\bigl(k^5(7-k)\bigr)+k\bigl(-4k^4\bigr)
      &&\text{[the two $3\times3$ minors]}\\
     &=2k^5(5-k)&&\text{[ring arithmetic].}
   \end{align*}
   Each $3\times3$ minor is its six-term determinant expansion. Thus C1 is a polynomial
   identity, certified by \path{Matrix.det_succ_row_zero} in
   \path{Mathlib/LinearAlgebra/Matrix/Determinant/Basic.lean}, with $D$ indexed by
   \texttt{Fin (3+1)}. For each minor use \path{Matrix.det_fin_three} in the same file,
   with index type \texttt{Fin 3}, then ring arithmetic.~\qed
  \end{proof}
  \qedstep
  \begin{proof}
   \pf\ Steps \stepref{210psd} and \stepref{210det} give $2k^5(5-k)\geq0$.
   Since $k\geq2$ by \stepref{210apply}, division by $2k^5>0$ gives $k\leq5$.~\qed
  \end{proof}
 \end{proof}

  \qedstep
  \begin{proof}\pf\ Step \stepref{210bound} is the claimed second proof.~\qed\end{proof}
 \end{proof}
 \qedstep
 \begin{proof}\pf\ Steps \stepref{210apply} and \stepref{210unique} give the assertions
  of the lemma; \stepref{210paperBound} records the paper's bound as well.~\qed\end{proof}
\end{proof}

\begin{cgss}{Remark}{2.11}\label{rem:2.11}
 \uses{cor:triangle-property}
 \emph{Side claim, not formalized.} We quote [CGSS]:
 \begin{quote}
 Similar methods may be used to prove the characterization of symplectic and orthogonal
 graphs over $\mathrm{GF}(2)$ by the triangle property, cf. [24, 21].
 \end{quote}
 We attach no Lean declaration to this side claim.
\end{cgss}

\begin{cgss}{Lemma}{2.12}\label{lem:2.12}
 \lean{LineSystem.Iso.ofOrthGraphIso}\lean{LineSystem.Hyp24.orthGraph_cases}
 \lean{LineSystem.IsIndecomposable.eq_pair_or_exists_hyp24}
 \uses{def:line-iso,def:indecomposable,hyp:2.4,lem:star-identities,lem:2.5,lem:2.6,lem:2.7,lem:2.9,lem:2.10,facts:inner}
 \assume{\begin{pfenum}
  \item $S$ is a nonempty indecomposable star-closed line system in a finite-dimensional
   real inner product space $V$.
  \item When a star is chosen, $\Ga$ and $X$ have the meanings of Hypothesis 2.4.
 \end{pfenum}}
 \prove{Either $S$ has one line, or it admits a pair $\Za,\Zb$ as in Hypothesis 2.4
  and, for every such choice, its graph $X$ is one of:
  the empty graph $\comp{K_m}$ on $m\geq0$ vertices; the friendship graph with
  $m\geq1$ triangles; the graphs of Lemma 2.10 with parameters
  $(9,4,1,2)$, $(15,6,1,3)$, or $(27,10,1,5)$.
  For any two instances of Hypothesis 2.4 in finite-dimensional real inner product
  spaces $V,W$, every isomorphism of their graphs on $\Ga$ extends to an isomorphism
  of their line systems, induced by a surjective linear isometry of their spans
  carrying the ordered stars to each other. This holds for every graph type in the list.}
\end{cgss}
\noindent\repair\ Nonemptiness excludes the empty line system, as required by
Implicit hypothesis 4 of \path{cgss/errata.md}.
\begin{proof}
 \pfsketch\ \repair\ We give all four strongly regular parameters in the exceptional
 cases to delimit the paper's uniqueness statement (Error 12). We will recover the signed root set
 from Lemmas 2.5 and 2.7.
 A graph isomorphism preserves the Gram matrix of the roots in
 $\{\Za,\Zb,\Zc\}\cup\Ga$, so it extends linearly to the spans and then carries
 the recovered sets to each other.

 \step{212star}{\repair\ Either $S$ consists of one antipodal pair or it admits Hypothesis 2.4.
  We include the one-line case omitted by the choice of a star in the paper
  (Implicit hypothesis 4).}
 \begin{proof}
  \pflet{Choose $\Zr\in\Phi$, using nonemptiness.}
  \step{212split}{If every root outside $\{\Zr,-\Zr\}$ is orthogonal to $\Zr$,
   then $\Phi=\{\Zr,-\Zr\}$.}
  \begin{proof}
   \pf\ Take $P=\{\Zr,-\Zr\}$ and $Q=\Phi\diff P$.
   The hypothesis and bilinearity make these parts orthogonal.
   Indecomposability and $\Zr\in P$ force $Q=\varnothing$.~\qed
  \end{proof}
  \step{212orient}{Otherwise there are $\Za,\Zb\in\Phi$ with $(\Za,\Zb)=-1$.}
  \begin{proof}
   \pf\ Choose $\Zs\notin\{\Zr,-\Zr\}$ with $(\Zr,\Zs)\ne0$.
   Equation \eqref{eq:endpoints} and the line-system condition give $\pm1$.
   Set $\Za=\Zr$, and set $\Zb=\Zs$ if the value is $-1$, or $\Zb=-\Zs$ if it is 1.
   Closure under negation supplies both choices.~\qed
  \end{proof}
  \qedstep
  \begin{proof}\pf\ The cases \stepref{212split} and \stepref{212orient} are exhaustive.~\qed\end{proof}
 \end{proof}

 \step{212recover}{In any instance of Hypothesis 2.4, put
  $T=\{\Za,\Zb,\Zc\}\cup\Ga$. Then $\Phi=\cR\cup(-\cR)$ for the set $\cR$
  in Lemma 2.7, and $\operatorname{span}(T)=\operatorname{span}(\Phi)$.}
 \begin{proof}
  \step{212formulas}{We have $\Phi=\cR\cup(-\cR)$, where
   \[
   \cR=\{\Za,\Zb,\Zc\}\cup\Ga\cup(\Zc+\Ga)\cup(-\Zb+\Ga)
   \cup\{\Zx-\Zy:\Zx,\Zy\in\Ga,\ (\Zx,\Zy)=1\}.
   \]}
  \begin{proof}\pf\ This is the recovery equality in the statement of Lemma 2.7,
   including both inclusions.~\qed\end{proof}
  \qedstep
  \begin{proof}
   \pf\ Step \stepref{212formulas} recovers the set using only $T$ and its inner products.
   Every displayed vector lies in $\operatorname{span}(T)$.
   Since $T\subseteq\Phi$, the two spans are equal.~\qed
  \end{proof}
 \end{proof}

 \pflet{For the extension assertion, fix any two instances of Hypothesis 2.4 in
  finite-dimensional real inner product spaces $V,W$ and a graph isomorphism
  $\psi:X\to Y$. Use tildes for the second root set and ordered star.
  For $x\in V(X)$, write $\Zs_x$ for the vector at $\psi(x)\in V(Y)$.}
 \step{212metric}{\new\ There is a surjective linear isometry between the two root spans,
  sending $\Za,\Zb,\Zc$ to $\widetilde{\Za},\widetilde{\Zb},\widetilde{\Zc}$
  and $\Zx$ to $\Zs_x$ for every $x\in V(X)$.
  We make the metric meaning of ``determined'' explicit.}
 \begin{proof}
  \pflet{Index two vector families $(\Zr_i),(\Zs_i)$ by the same finite disjoint union
   $I=\{1,2,3\}\sqcup V(X)$.
   The first three values are the ordered stars; the other values are
   $\Zr_x=\Zx$ and the vectors $\Zs_x$ just defined.
   Put $R(\alpha)=\sum_i\alpha(i)\Zr_i$ and $R'(\alpha)=\sum_i\alpha(i)\Zs_i$.
   These are linear maps from $I\to\re$ to $V,W$, respectively, by
   \path{Fintype.linearCombination} in
   \path{Mathlib/LinearAlgebra/Finsupp/LinearCombination.lean}.}
  \step{212gram}{$(\Zr_i,\Zr_j)=(\Zs_i,\Zs_j)$ for all $i,j\in I$.}
  \begin{proof}
   \step{212stargram}{Products within the first three positions are 2 on the diagonal
    and $-1$ off it; products with a vertex of $X$ are $0,1,-1$ in that order.}
   \begin{proof}\pf\ These are the star identities and the defining products of $\Ga$
    in each system.~\qed\end{proof}
   \step{212graphgram}{Products within the graph positions are 2 on the diagonal,
    0 on edges, and 1 on nonedges.}
   \begin{proof}
    \pf\ Use Lemma 2.6 and the orthogonality convention; the bijection $\psi$ preserves
    equality, adjacency, and nonadjacency.~\qed
   \end{proof}
   \qedstep
   \begin{proof}\pf\ Steps \stepref{212stargram} and \stepref{212graphgram}, with symmetry, cover all pairs of positions.~\qed\end{proof}
  \end{proof}
  \step{212relations}{$\ker R\subseteq\ker R'$: for every $\alpha:I\to\re$,
   $R(\alpha)=0$ implies $R'(\alpha)=0$.}
  \begin{proof}
   \step{212norms}{The two sums have the same squared norm.}
   \begin{proof}
    \pf\ Expanding by bilinearity gives
    \[
    \sum_{i,j}\alpha(i)\alpha(j)(\Zr_i,\Zr_j)
     =\sum_{i,j}\alpha(i)\alpha(j)(\Zs_i,\Zs_j)
    \]
    by \stepref{212gram}.~\qed
   \end{proof}
   \qedstep
   \begin{proof}
    \pf\ If the first sum is zero, \stepref{212norms} makes the second sum have squared
    norm zero. The zero-norm implication in the elementary identities makes that sum zero.~\qed
   \end{proof}
  \end{proof}
  \step{212linear}{There is a surjective linear map $F$ of the two spans such that
   \[F\left(\sum_i\alpha(i)\Zr_i\right)=\sum_i\alpha(i)\Zs_i.\]}
  \begin{proof}
   \step{212ranges}{The ranges of $R,R'$ are the spans of $(\Zr_i),(\Zs_i)$, respectively.}
   \begin{proof}
    \raggedright\pf\ Apply \path{Fintype.range_linearCombination} in
    \path{Mathlib/LinearAlgebra/Finsupp/LinearCombination.lean} to each finite family.~\qed
   \end{proof}
   \pflet{Let $R'_0$ be $R'$ with codomain restricted to $\operatorname{range}R'$,
    using \path{LinearMap.rangeRestrict} in
    \path{Mathlib/Algebra/Module/Submodule/Range.lean}.}
   \step{212kernel}{$\ker R\subseteq\ker R'_0$.}
   \begin{proof}
    \pf\ A vector in the subspace $\operatorname{range}R'$ is zero exactly when its
    underlying vector is zero. Thus \stepref{212relations} gives the inclusion.~\qed
   \end{proof}
   \step{212quotient}{There is a linear map
    $\overline R':(I\to\re)/\ker R\to\operatorname{range}R'$
    satisfying $\overline R'([\alpha])=R'_0(\alpha)$.}
   \begin{proof}
    \pf\ Apply \path{Submodule.liftQ} to $R'_0$ and \stepref{212kernel};
    \path{Submodule.liftQ_apply} gives the formula. Both declarations are in
    \path{Mathlib/LinearAlgebra/Quotient/Basic.lean}.~\qed
   \end{proof}
   \step{212rangemap}{There is a linear map
    $F_0:\operatorname{range}R\to\operatorname{range}R'$
    with $F_0(R(\alpha))=R'(\alpha)$ for all $\alpha$.}
   \begin{proof}
    \pflet{Let $E:(I\to\re)/\ker R\to\operatorname{range}R$ be
     \path{LinearMap.quotKerEquivRange} for $R$, in
     \path{Mathlib/LinearAlgebra/Isomorphisms.lean}, and put $F_0=\overline R'\circ E^{-1}$.}
    \step{212compose}{$F_0$ is linear.}
    \begin{proof}\pf\ It is a composite of the linear map in \stepref{212quotient}
     and the inverse of a linear equivalence.~\qed\end{proof}
    \step{212oncombinations}{$F_0(R(\alpha))=R'(\alpha)$.}
    \begin{proof}
     \pf\ The lemma \path{LinearMap.quotKerEquivRange_symm_apply_image} in
     \path{Mathlib/LinearAlgebra/Isomorphisms.lean} gives $E^{-1}(R(\alpha))=[\alpha]$.
     Apply the formula in \stepref{212quotient}.~\qed
    \end{proof}
    \qedstep
    \begin{proof}\pf\ Steps \stepref{212compose} and \stepref{212oncombinations}
     prove both assertions.~\qed\end{proof}
   \end{proof}
   \step{212linearonto}{$F_0$ is surjective.}
   \begin{proof}\pf\ Every member of $\operatorname{range}R'$ has the form $R'(\alpha)$;
    by \stepref{212rangemap}, $R(\alpha)$ is a preimage.~\qed\end{proof}
   \step{212generatorvalues}{$F_0(\Zr_i)=\Zs_i$ for every $i$.}
   \begin{proof}
    \pf\ In \stepref{212rangemap}, take the coefficient function which is 1 at $i$
    and zero elsewhere. The two finite sums are $\Zr_i,\Zs_i$.~\qed
   \end{proof}
   \qedstep
   \begin{proof}\pf\ Identify the ranges with the spans by \stepref{212ranges}.
    The map $F_0$ in \stepref{212rangemap}, with surjectivity from
    \stepref{212linearonto}, is the required $F$ and has the values in
    \stepref{212generatorvalues}.~\qed\end{proof}
  \end{proof}
  \pflet{Choose a surjective linear map $F$ with the formula in \stepref{212linear}.}
  \step{212isometry}{The map $F$ preserves inner products and is injective.}
  \begin{proof}
   \step{212inner}{For vectors $\Zx,\Zy$ in the first span,
    $(F(\Zx),F(\Zy))=(\Zx,\Zy)$.}
   \begin{proof}
    \pf\ By the finite-span lemma, write $\Zx,\Zy$ as combinations with coefficients
    $\alpha,\beta$. The formula in \stepref{212linear} gives their images, and bilinearity gives
    \[
     \sum_{i,j}\alpha(i)\beta(j)(\Zr_i,\Zr_j)
       =\sum_{i,j}\alpha(i)\beta(j)(\Zs_i,\Zs_j)
    \]
    by \stepref{212gram}.~\qed
   \end{proof}
   \step{212injective}{$F$ is injective.}
   \begin{proof}
    \pf\ If $F(\Zx)=F(\Zy)$, linearity gives $F(\Zx-\Zy)=0$.
    Step \stepref{212inner} gives $(\Zx-\Zy,\Zx-\Zy)=0$.
    The zero-norm identity then gives $\Zx=\Zy$.~\qed
   \end{proof}
   \qedstep
   \begin{proof}\pf\ Combine \stepref{212inner} and \stepref{212injective}.~\qed\end{proof}
  \end{proof}
  \qedstep
  \begin{proof}
   \pf\ Steps \stepref{212linear} and \stepref{212isometry} give a linear isometry
   onto the second span. Step \stepref{212recover} identifies these spans with the root spans.
   Taking one coefficient equal to 1 and all others zero in \stepref{212linear}
   gives its prescribed values on the two indexed families.
   This is the equal-Gram argument LA2 from the prerequisites, proved here for the needed families.~\qed
  \end{proof}
 \end{proof}

 \pflet{Choose an isometry $F$ with the values and surjectivity in \stepref{212metric}.}
 \step{212rootiso}{The isometry $F$ in \stepref{212metric} carries the first root set
  bijectively onto the second.}
 \begin{proof}
  \step{212image}{$F(\Phi)\subseteq\widetilde\Phi$.}
  \begin{proof}
   \step{212imagea}{$F(\Za)=\widetilde{\Za}\in\widetilde\Phi$.}
   \begin{proof}\pf\ Use \stepref{212metric} and the second instance of Hypothesis 2.4.~\qed\end{proof}
   \step{212imageb}{$F(\Zb)=\widetilde{\Zb}\in\widetilde\Phi$.}
   \begin{proof}\pf\ Use \stepref{212metric} and the second instance of Hypothesis 2.4.~\qed\end{proof}
   \step{212imagec}{$F(\Zc)=\widetilde{\Zc}\in\widetilde\Phi$.}
   \begin{proof}\pf\ Use \stepref{212metric} and the star identities in the second instance.~\qed\end{proof}
   \step{212imagegamma}{For $x\in V(X)$, $F(\Zx)=\Zs_x\in\widetilde\Ga$.}
   \begin{proof}\pf\ Step \stepref{212metric} gives the value; $\Zs_x$ is the
    vector at $\psi(x)$, a vertex of the second graph.~\qed\end{proof}
   \step{212imageplus}{For $\Zx\in\Ga$, $F(\Zc+\Zx)=\widetilde{\Zc}+\Zs_x\in\widetilde\Phi$.}
   \begin{proof}\pf\ Linearity and \stepref{212metric} give the equality.
    Lemma 2.6 in the second instance, or the reverse inclusion of its recovery
    equality in Lemma 2.7, gives membership.~\qed\end{proof}
   \step{212imageminus}{For $\Zx\in\Ga$, $F(-\Zb+\Zx)=-\widetilde{\Zb}+\Zs_x\in\widetilde\Phi$.}
   \begin{proof}\pf\ Linearity and \stepref{212metric} give the equality.
    Apply the second translation in Lemma 2.6 in the target, equivalently the
    reverse inclusion of its recovery equality in Lemma 2.7.~\qed\end{proof}
   \step{212imagedifference}{If $\Zx,\Zy\in\Ga$ and $(\Zx,\Zy)=1$,
    then $F(\Zx-\Zy)=\Zs_x-\Zs_y\in\widetilde\Phi$.}
   \begin{proof}
    \step{212differenceproduct}{$(\Zs_x,\Zs_y)=1$.}
    \begin{proof}\pf\ The isometry in \stepref{212metric} preserves the given product.~\qed\end{proof}
    \qedstep
    \begin{proof}\pf\ Linearity gives the image formula. The reverse inclusion
     of the second recovery equality in Lemma 2.7 applies to
     \stepref{212differenceproduct}.~\qed\end{proof}
   \end{proof}
   \step{212imagesigns}{If a displayed expression $\Zr$ maps into $\widetilde\Phi$,
    so does $-\Zr$.}
   \begin{proof}\pf\ Linearity gives $F(-\Zr)=-F(\Zr)$ and
    $\widetilde\Phi=-\widetilde\Phi$.~\qed\end{proof}
   \qedstep
   \begin{proof}
    \pf\ The forward inclusion of the first recovery equality in Lemma 2.7
    expresses every root by one of these seven kinds and a sign.
    Steps \stepref{212imagea} through \stepref{212imagedifference}
    handle the expressions, and \stepref{212imagesigns} handles negation.~\qed
   \end{proof}
  \end{proof}
  \step{212onto}{$\widetilde\Phi\subseteq F(\Phi)$.}
  \begin{proof}
   \step{212ontostar}{Each root of the second ordered star has its corresponding
    first-star root as a preimage in $\Phi$.}
   \begin{proof}\pf\ Apply the prescribed star values in \stepref{212metric};
    Hypothesis 2.4 and the star identities put all three preimages in $\Phi$.~\qed\end{proof}
   \step{212ontogamma}{Each $\Zs\in\widetilde\Ga$ has a preimage $\Zx\in\Ga$.}
   \begin{proof}\pf\ Take the inverse image of its vertex under the bijection $\psi$
    and use \stepref{212metric}.~\qed\end{proof}
   \step{212ontoplus}{Each $\widetilde{\Zc}+\Zs$ with $\Zs\in\widetilde\Ga$
    has a preimage in $\Phi$.}
   \begin{proof}\pf\ Choose $\Zx$ by \stepref{212ontogamma}.
    The reverse inclusion of the first recovery equality in Lemma 2.7 puts
    $\Zc+\Zx$ in $\Phi$, and linearity gives its image.~\qed\end{proof}
   \step{212ontominus}{Each $-\widetilde{\Zb}+\Zs$ with $\Zs\in\widetilde\Ga$
    has a preimage in $\Phi$.}
   \begin{proof}\pf\ Choose $\Zx$ by \stepref{212ontogamma}.
    The reverse inclusion of the first recovery equality in Lemma 2.7 puts
    $-\Zb+\Zx$ in $\Phi$, and linearity gives its image.~\qed\end{proof}
   \step{212ontodifference}{Each $\Zs-\Zt$ with $\Zs,\Zt\in\widetilde\Ga$
    and $(\Zs,\Zt)=1$ has a preimage in $\Phi$.}
   \begin{proof}
    \pf\ Choose preimages $\Zx,\Zy$ by \stepref{212ontogamma}.
    The isometry in \stepref{212metric} gives $(\Zx,\Zy)=1$.
    The reverse inclusion of the first recovery equality in Lemma 2.7 puts
    $\Zx-\Zy$ in $\Phi$, and linearity gives its image.~\qed
   \end{proof}
   \step{212ontosigns}{The negative of a target with a preimage in $\Phi$
    also has a preimage in $\Phi$.}
   \begin{proof}\pf\ Negate the preimage, using $\Phi=-\Phi$ and linearity.~\qed\end{proof}
   \qedstep
   \begin{proof}
    \pf\ The forward inclusion of the second recovery equality in Lemma 2.7
    lists exactly the expressions and signs covered by \stepref{212ontostar}
    through \stepref{212ontosigns}.~\qed
   \end{proof}
  \end{proof}
  \qedstep
  \begin{proof}
   \pf\ Injectivity in \stepref{212metric}, \stepref{212image}, and \stepref{212onto}
   make the restriction of $F$ a bijection preserving inner products, as required
   by the definition of a line-system isomorphism.~\qed
  \end{proof}
 \end{proof}

 \step{212graphs}{Every graph arising from a chosen star belongs to the stated list.}
 \begin{proof}
  \step{212empty}{\case{$X$ has no edges.}}
  \begin{proof}\pf\ Its finite vertex set has some cardinality $m\geq0$, giving the first type.~\qed\end{proof}
  \step{212universal}{\case{$X$ has an edge and has a universal vertex.}}
  \begin{proof}
   \pf\ Lemma 2.9 gives a friendship graph with $m$ triangles.
   Since $X$ has an edge, $m\ne0$, so $m\geq1$.~\qed
  \end{proof}
  \step{212remaining}{\case{$X$ has an edge and has no universal vertex.}}
  \begin{proof}
   \pf\ A complete graph with an edge has universal vertices, so this graph is noncomplete.
   Thus Lemma 2.10 applies; its statement gives the displayed strongly regular parameters.~\qed
  \end{proof}
  \qedstep
  \begin{proof}\pf\ These three cases exhaust existence of an edge and of a universal vertex.~\qed\end{proof}
 \end{proof}
 \qedstep
 \begin{proof}
  \pf\ Step \stepref{212star} separates the one-line case; \stepref{212graphs} lists
  every graph in the other case. Steps \stepref{212metric} and \stepref{212rootiso}
  give the claimed uniqueness and its realization by a linear isometry of spans.~\qed
 \end{proof}
\end{proof}

\paragraph{Inputs from Section 3.}\label{input:sec3-models}
\repair\ The line counts and dimensions in Theorem 2.13 require the models of Section 3 (Gap 4 of \path{cgss/errata.md}).
We will use the following facts from Definitions 3.1 to 3.3, Lemma 3.4, and computations C2
and C4. We record them as deferred dependencies, not as conclusions of Section 2.
The labels \path{def:3.1}, \path{def:3.2}, \path{def:3.3}, \path{lem:3.4},
\path{comp:C2} and \path{comp:C4} are those of Section 3, which proves these facts.
Each model is a nonempty indecomposable star-closed line system in its real span. The model
$\typ{A}_1$ consists of one antipodal pair and has rank 1. Each remaining model has a specified
ordered star for which the following data hold. Here rank means dimension of the real span,
and $\delta$ means $|\Delta|/2$, the number of lines in $\Delta$.
\[
\begin{array}{c|c|c|c|c}
 \text{model}&\text{rank}&\text{number of lines}&\text{graph on }\Ga&\delta\\\hline
 \typ{A}_\ell\ (\ell\geq2)&\ell&\ell(\ell+1)/2&\comp{K_{\ell-2}}&\text{not used}\\
 \typ{D}_\ell\ (\ell\geq4)&\ell&\ell(\ell-1)&\text{friendship, }\ell-3\text{ triangles}&\text{not used}\\
 \typ{E}_6&6&36&(9,4,1,2)&6\\
 \typ{E}_7&7&63&(15,6,1,3)&15\\
 \typ{E}_8&8&120&(27,10,1,5)&36
\end{array}
\]
In the exceptional rows the tuple gives strongly regular parameters; the graph is the
orthogonality graph, so it has degree 4, 6, or 10, respectively. We take its uniqueness from
Lemma 2.10, not from the model verification. C2 gives the line counts and C4 verifies
star-closure, indecomposability and the marked-star data; spanning sets in the explicit
models give the displayed ranks. These verifications must not use Theorem 2.13 or Theorem 3.5.
For Definition 3.3 we use fixed representatives of a root and a star in $\typ{E}_8$.
The fixed-root construction of $\typ{E}_7$ and the fixed-star construction of $\typ{E}_6$
suffice here. Prerequisite LN5 is needed to compare a fixed-root construction with the
phrase ``any one of its lines'' in Definition 3.3, but is not used in Theorem 2.13.

\begin{cgss}{Theorem}{2.13}\label{thm:2.13}
 \lean{LineSystem.five_possibilities}\lean{LineSystem.five_possibilities'}
 \lean{LineSystem.five_possibilities_forall_hyp24}
 \lean{LineSystem.Iso.exists_ambient_isometry_of_same_model}
 \lean{LineSystem.Hyp24.six_lines_of_iso_friendship_zero}
 \uses{lem:2.12,lem:2.10,lem:2.9,def:line-iso,facts:inner,input:sec3-models,def:3.1,def:3.2,def:3.3,lem:3.4,comp:C2,comp:C4}
 \assume{\begin{pfenum}
  \item $V$ is a finite-dimensional real inner product space of dimension $\ell$.
  \item $S$ is a nonempty indecomposable star-closed line system in $V$, and
   $\operatorname{span}(\Phi)=V$.
 \end{pfenum}}
 \prove{The possibilities for $S$, up to isomorphism, are the five cases below.
 For every case with more than one line, there exist $\Za,\Zb\in\Phi$ with
 $(\Za,\Zb)=-1$ for which the stated data on $\Ga$ and $\Delta$ hold.
 \repair\ We exclude the empty system and separate the one-line case
 (Implicit hypothesis 4 of \path{cgss/errata.md}).
 \begin{pfenum}
  \item $\ell(\ell+1)/2$ lines, with the empty graph $\comp{K_{\ell-2}}$ on $\Ga$ when $\ell\geq2$;
   for $\ell=1$ there is one line and no chosen star;
  \item $\ell(\ell-1)$ lines for $\ell\geq4$, with the friendship graph of $\ell-3$ triangles on $\Ga$;
  \item 120 lines in dimension 8, with $|\Ga|=27$ and $|\Delta|/2=36$;
  \item 63 lines in dimension 7, with $|\Ga|=15$ and $|\Delta|/2=15$;
  \item 36 lines in dimension 6, with $|\Ga|=9$ and $|\Delta|/2=6$.
 \end{pfenum}
 Each case is unique up to an orthogonal transformation when the ambient space is $\re^\ell$.
 The zero-triangle case of Lemma 2.9 has $\ell=3$ and six lines; it is already in case 1
 and also satisfies the formula $\ell(\ell-1)=6$.}
 The Lean statement gives the data on $\Ga$ and $\Delta$ for some star, while in the paper they
 belong to the star fixed after Hypothesis 2.4, which is arbitrary. Every star gives the same
 data: \path{LineSystem.five_possibilities_forall_hyp24} is the theorem for every star
 $(\Za,\Zb)$ of Hypothesis 2.4 whose system spans $V$, with the graph on $\Ga$ and the counts of
 $\Ga$ and $\Delta$ taken at that star. It leaves out the one-line case, which has no star.
\end{cgss}
\begin{proof}
 \pfsketch\ We will match the graph from Lemma 2.12 with a specified star in one of the
 models. Its graph isomorphism gives a linear isometry of spans, so it transfers both
 rank and number of lines. Thus the numerical conclusions come from Section 3.

 \step{213one}{If $S$ has one line, then $\ell=1$ and a linear isometry of spans
  carries it onto $\typ{A}_1$.}
 \begin{proof}
  \step{213rankone}{The span of $\Phi=\{\Zr,-\Zr\}$ is one-dimensional.}
  \begin{proof}
   \pf\ Squared norm 2 gives $\Zr\ne0$. Thus the singleton $\{\Zr\}$ is a basis
   of this span, which equals $V$ by \asm[2]{0}.~\qed
  \end{proof}
  \step{213onemap}{For either root $\Zs$ of $\typ{A}_1$, the rule
   $F(\alpha\Zr)=\alpha\Zs$ defines a surjective linear isometry of spans.}
  \begin{proof}
   \pf\ Step \stepref{213rankone} gives a unique coefficient $\alpha$ for each vector.
   Addition and scaling of this coefficient make $F$ linear and surjective.
   For any real $\alpha,\beta$, both inner products are $2\alpha\beta$ by bilinearity.
   Thus $F$ preserves inner products.~\qed
  \end{proof}
  \qedstep
  \begin{proof}
   \pf\ The map in \stepref{213onemap} sends $\{\Zr,-\Zr\}$ bijectively onto
   $\{\Zs,-\Zs\}$. The line count is $1=1(1+1)/2$, and
   \stepref{213rankone} gives $\ell=1$.~\qed
  \end{proof}
 \end{proof}

 \step{213match}{If $S$ has more than one line, its chosen-star graph is isomorphic
  to the specified graph of one model in the input table.}
 \begin{proof}
  \pflet{Choose a star using Lemma 2.12 and let $X$ be its graph.}
  \step{213A}{\case{$X=\comp{K_m}$.}}
  \begin{proof}
   \pf\ The specified star of $\typ{A}_{m+2}$ has the empty graph on $(m+2)-2=m$ vertices
   by the Section 3 inputs. A bijection between the finite vertex sets preserves
   adjacency since both graphs have no edges.~\qed
  \end{proof}
  \step{213D}{\case{$X$ is the friendship graph with $m\geq1$ triangles.}}
  \begin{proof}
   \step{213Dsize}{The specified graph in $\typ{D}_{m+3}$ is the friendship graph with $m$ triangles.}
   \begin{proof}\pf\ Use $m+3\geq4$ and the $\typ{D}_\ell$ row of the Section 3 inputs.~\qed\end{proof}
   \qedstep
   \begin{proof}
    \pf\ Send a chosen common vertex to the model's common vertex, choose a bijection
    of the $m$ pairs of remaining vertices, and a bijection within each pair.
    Lemma 2.9 lists exactly the edges preserved by this map, so it is a graph isomorphism.~\qed
   \end{proof}
  \end{proof}
  \step{213E}{\case{$|V(X)|\in\{9,15,27\}$ as in Lemma 2.10.}}
  \begin{proof}
   \step{213Echoose}{Choose $\typ{E}_6,\typ{E}_7,\typ{E}_8$ according as the order is $9,15,27$.
    Its specified graph satisfies the hypotheses of Lemma 2.10.}
   \begin{proof}
    \pf\ The respective orders and degrees in the input table are $(9,4),(15,6),(27,10)$.
    Each degree is positive and smaller than order minus 1, so each graph has an edge, is
    noncomplete, and has no universal vertex. A graph of Lemma 2.9 has a universal
    vertex, so that type is excluded.~\qed
   \end{proof}
   \qedstep
   \begin{proof}
    \pf\ The chosen model is a star-closed indecomposable line system and has the same
    graph order as $X$. Apply the uniqueness assertion of Lemma 2.10.~\qed
   \end{proof}
  \end{proof}
  \qedstep
  \begin{proof}\pf\ Lemma 2.12 makes \stepref{213A}, \stepref{213D}, and \stepref{213E} exhaustive.~\qed\end{proof}
 \end{proof}

 \step{213transfer}{For a model matched in \stepref{213match}, $\ell$ equals its rank and
  $S$ has its number of lines. The marked sets $\Delta$ also have the same number of lines.}
 \begin{proof}
  \step{213spaniso}{The graph isomorphism extends to a surjective linear isometry $F$
   between the spans, carrying the ordered stars to each other and mapping root sets bijectively.}
  \begin{proof}\pf\ Apply the reconstruction and equal-Gram argument of Lemma 2.12.~\qed\end{proof}
  \step{213dimension}{The two spans have equal dimensions.}
  \begin{proof}
   \pf\ A linear isometry is injective; together with surjectivity in
   \stepref{213spaniso}, it sends a basis bijectively to a basis.
   The first span is $V$ by \asm[2]{0}.~\qed
  \end{proof}
  \step{213lines}{The root bijection induces a bijection of antipodal pairs.}
  \begin{proof}
   \pf\ Linearity gives $F(-\Zr)=-F(\Zr)$; injectivity prevents two different
   antipodal pairs from merging, and surjectivity on roots covers every target pair.~\qed
  \end{proof}
  \step{213delta}{The root bijection sends $\Delta$ exactly to the model's marked $\Delta$.}
  \begin{proof}
   \pf\ The definition of $\Delta$ asks for three zero inner products with the ordered
   star. Both $F$ and its inverse preserve these products and the ordered star by
   \stepref{213spaniso}. Thus membership is preserved in both directions.~\qed
  \end{proof}
  \qedstep
  \begin{proof}
   \pf\ Step \stepref{213dimension} transfers rank, \stepref{213lines} transfers line count,
   and \stepref{213delta} with \stepref{213lines} transfers the line count in $\Delta$.~\qed
  \end{proof}
 \end{proof}

 \step{213numbers}{The transferred data give exactly the five numerical cases in the statement.}
 \begin{proof}
  \step{213classical}{The empty-graph case gives $\ell=m+2$ and $\ell(\ell+1)/2$ lines;
   the friendship case with $m\geq1$ triangles gives $\ell=m+3$ and $\ell(\ell-1)$ lines.}
  \begin{proof}
   \pf\ Substitute the ranks of $\typ{A}_{m+2},\typ{D}_{m+3}$ from the input table into
   \stepref{213transfer}, then use their line-count formulas from Lemma 3.4 and C2.~\qed
  \end{proof}
  \step{213exceptional}{The orders $27,15,9$ give respectively the data
   \[(\ell,|\Phi|/2,|\Delta|/2)=(8,120,36),(7,63,15),(6,36,6).\]}
  \begin{proof}
   \pf\ Use the $\typ{E}_8,\typ{E}_7,\typ{E}_6$ rows of the Section 3 inputs in \stepref{213transfer}.
   \repair\ We correct Error 6: we write $|\Delta|/2$ because the signed set $\Delta$ has respectively
   $72,30,12$ vectors. The paper's $|\Delta|=36,15,6$ counts lines.~\qed
  \end{proof}
  \step{213overlap}{\new\ The zero-triangle case of Lemma 2.9 has six lines in dimension 3; we make this overlap explicit.}
  \begin{proof}
   \pf\ Its graph is a single vertex, the empty graph $\comp{K_1}$.
   The first formula in \stepref{213classical} gives $\ell=3$ and $3\cdot4/2=6=3\cdot2$.~\qed
  \end{proof}
  \qedstep
  \begin{proof}
   \pf\ Combine \stepref{213one}, \stepref{213classical}, \stepref{213exceptional},
   and \stepref{213overlap}.~\qed
  \end{proof}
 \end{proof}

 \step{213unique}{Each numerical case, with its stated graph type, is unique up to
  an orthogonal transformation in $\re^\ell$.}
 \begin{proof}
  \step{213samemodel}{Two systems in the same case are isomorphic to the same model.}
  \begin{proof}
   \pf\ In the first family the model rank is $\ell$, so the model is $\typ{A}_\ell$;
   in the second it is $\typ{D}_\ell$; the exceptional rows specify $\typ{E}_8,\typ{E}_7,\typ{E}_6$ individually.
   Use \stepref{213one} for $\ell=1$ and \stepref{213match} otherwise.~\qed
  \end{proof}
  \qedstep
  \begin{proof}
   \pf\ Compose the span isometry to that model with the inverse for the other system.
   Both root spans are the entire $\re^\ell$ by \asm[2]{0}. Thus the composite is a
   surjective linear isometry of $\re^\ell$, an orthogonal transformation, sending one
   line system onto the other.~\qed
  \end{proof}
 \end{proof}
 \qedstep
 \begin{proof}
  \pf\ Step \stepref{213numbers} gives the alternatives and \stepref{213unique} gives
  uniqueness. The asserted dimensions and counts depend on the explicit Section 3 inputs
  recorded above; no classification result of Section 3 was used.~\qed
 \end{proof}
\end{proof}
